\documentclass[aps,pra,reprint,superscriptaddress,longbibliography]{revtex4-2}
\usepackage[T1]{fontenc}
% A partial Latin Modern cache can lack the five-point roman metric.
% Full installations keep their normal font definitions unchanged.
\IfFileExists{ec-lmr5.tfm}{}{%
 \IfFileExists{t1lmr.fd}{%
  \input{t1lmr.fd}%
  \DeclareFontShape{T1}{lmr}{m}{n}
   {<-6.5>ec-lmr6 <6.5-7.5>ec-lmr7 <7.5-8.5>ec-lmr8
    <8.5-9.5>ec-lmr9 <9.5-11>ec-lmr10 <11-15>ec-lmr12
    <15->ec-lmr17}{}%
 }{}%
}
\usepackage{amsmath,amssymb,bm,booktabs,graphicx}
\usepackage[colorlinks=true,linkcolor=blue,citecolor=blue,urlcolor=blue]{hyperref}
\newcommand{\Tr}{\operatorname{Tr}}
\newcommand{\dd}{\mathrm d}
\newcommand{\ee}{\mathrm e}
\newcommand{\id}{\mathbb I}
\newcommand{\supp}{\operatorname{supp}}
\newcommand{\Var}{\operatorname{Var}}
\newcommand{\sinc}{\operatorname{sinc}}
\newcommand{\norm}[1]{\left\lVert#1\right\rVert}
\newcommand{\ket}[1]{\lvert#1\rangle}
\newcommand{\bra}[1]{\langle#1\rvert}
\newcommand{\ip}[2]{\langle#1,#2\rangle}
\newcommand{\one}{\mathbf 1}
\newcommand{\Srel}{\mathsf D}
\begin{document}
\title{Boundary coherence and work fluctuations under quantum confinement control}
\author{R. Arun Chandru}
\email{arun.chandru@gmail.com}
\date{\today}
\begin{abstract}
Coherent boundary interference can purify the retained part of an
increasingly mixed Gibbs state, while its work statistics probe
inequivalent spectral sectors. We establish this connection for a confined
particle subject to a coherently controlled potential wall. A common
boundary source governs the conditional state, while its positive
spectral measure determines the leading population--coherence
allocation of dissipated work and the Ramsey response. Resonant
interference produces an asymptotically
pure retained state with an explicitly paid success probability.
For a sharp wall, a discrete near-elastic sector contributes a
logarithm to absolute work that is absent from the quadratic-weight
continuum limit. Its contribution makes the ratio of absolute to net
work diverge on specified thermal paths. Smooth-wall thermodynamic
limits, a finite smooth-device certificate with timing and readout
errors, and a finite-count reconstruction scheme connect these results
to operational resources.
\end{abstract}
\maketitle

\section{Introduction}

Coherent control and thermal averaging interrogate different aspects of
a quantum boundary. Thermal occupation counts the states available in
an extended well. A local change of the wall instead couples to the
small part of each wave function that penetrates the controlled region.
It is therefore possible for many orthogonal occupied states to have
nearly parallel responses to the same local perturbation. What remains
to be established is whether this alignment survives the complete
Gibbs ensemble, how much a retained state costs, and which work
statistics can be inferred from its interferometric signature.

The characteristic function of quantum work can be measured by a
Ramsey interferometer, without carrying out the two projective energy
measurements in each run \cite{Dorner2013,Mazzola2013,Batalhao2014}.
Conditional purification through repeated ancillary measurements is
also well established \cite{Nakazato2003}, including thermal-state
filtering by Ramsey postselection \cite{Oyarce2019} and engineered
ancilla-mediated logical-state protocols \cite{Pushpan2025}.
These observations motivate
the experiment considered here, but do not settle its continuum limit.
In particular, a vanishing unnormalized error need not remain small
after conditioning on a rare outcome. Nor does convergence of purity
to one imply convergence of entropy in an unbounded Hilbert space.
The retained state, its full success probability, and its spectral
tails must be controlled together.

Another distinction concerns the information contained in work.
The decomposition of irreversible entropy production into population
and coherence contributions is established in quantum thermodynamics
\cite{Francica2019}. In an extended continuum system, the thermal scale
can compete with the exponentially small Gibbs weight at a fixed final
energy. The boundary contribution must then be followed through the
moving crossover rather than expanded uniformly in the transition
probability. Finite-dimensional high-temperature quench expansions
\cite{Varizi2022} address a different limit: here the barrier is
thermally high and the number of populated low modes diverges.
We obtain an explicit capped spectral moment for this
competition, including the thermally activated population above the
initial barrier.

The same system also reveals a limitation of inferring singular work
functionals from a regular interferometric limit. A quadratic work
weight gives a finite limiting measure. An absolute-value weight is
more sensitive to transitions between nearby levels. We calculate
their accumulated contribution directly in the finite well. It is
logarithmic in the number of thermally occupied modes and survives
after the regular continuum contribution has been extracted. Thus
fixed-time convergence of a Ramsey response does not justify
interchanging the continuum limit with an unbounded observation-time
integral. Local-quench work edges and scattering-controlled Ramsey
responses already occur in critical and Fermi systems
\cite{Smacchia2012,Knap2012,Sindona2013}. Scattering-induced logarithms
are also central to orthogonality catastrophe \cite{Kuttler2014}.
The present calculation instead concerns a single-particle Gibbs
ensemble, a finite wall-height ratio, and a gap-weighted thermal
harmonic sum with an additive remainder. Its finite-amplitude quench
and singular thermal limit distinguish it from weak-amplitude work
response \cite{Guarnieri2024}.

In an earlier work~\cite{ChandruRareWork}, we established finite-wall boundary-source
laws and rare work costs at fixed inverse temperature in box units.
Here that inverse temperature tends to zero, and the number of thermally
occupied well modes diverges jointly with the barrier scale. The resulting
coherence is leading order, the relevant thermal crossover is retained,
and the nearby-level work sector requires a separate calculation.
The physical instrument remains a pair of controlled propagators;
no final motional-energy projection or commuting surrogate filter is
inserted. The analysis yields four connected results. It identifies
the common boundary spectral measure, determines its thermodynamic
allocation, resolves the singular absolute-work law for a sharp wall,
and gives finite operational statements for purification and
reconstruction. The complete derivations below make explicit where
a smooth-profile result applies and where the sharp-wall calculation
uses its exact matching conditions.

\section{A controlled boundary and its physical scales}
\label{sec:model}

Let a particle occupy the interval $0<y<2$, with Dirichlet boundary
conditions. Initially and on the controlled branch its dimensionless
Hamiltonians are
\begin{equation}
 A_r=-r^2\partial_y^2+F((y-1)/r),\qquad
 B_r=A_r+3F((y-1)/r).                         \label{eq:ham}
\end{equation}
The sharp wall is $F(x)=\one_{(0,\infty)}(x)$. The smooth asymptotic
family consists of a fixed nondecreasing smooth $F$ with $0\le F\le1$,
$F=0$ on $x\le-d$ and $F=1$ on $x\ge d$, for a fixed $d>0$.
Constants may depend on this profile. A separate finite example will
use an error-function wall with noncompact, exponentially small tails.

The prepared state is the full Gibbs state
\begin{equation}
 \rho_{r,s}=Z_1^{-1}\ee^{-sA_r},\qquad
 Z_h=\Tr\ee^{-s[-r^2\partial_y^2+hF((y-1)/r)]}.
                                                               \label{eq:gibbs}
\end{equation}
All initial eigenstates are included. A pure ancillary two-level
system coherently selects the two propagators. Its dark outcome
retains the motional operator
\begin{equation}
 \begin{gathered}
 K_t=\frac{\ee^{-itB_r}-\ee^{-itA_r}}2,\quad\\
 p_{r,s}(t)=\Tr K_t\rho_{r,s}K_t^*,\quad\\
 \sigma_t=K_t\rho_{r,s}K_t^*/p_{r,s}(t).          \label{eq:kraus}
 \end{gathered}
\end{equation}
The complementary Kraus operator is the sum of the same propagators
divided by two. Their squared operators add to the identity.

To restore units, let the physical coordinate be $L_0y$ and put
$E_0=\hbar^2/(2mL_0^2)$. Then
\begin{equation}
 H_{A,B}=\frac{E_0}{r^2}(A_r,B_r),\qquad
 \beta_{\rm phys}=\frac{s r^2}{E_0},\qquad
 t_{\rm phys}=\frac{\hbar r^2}{E_0}t.             \label{eq:units}
\end{equation}
The barrier height grows as $r^{-2}$, while a fixed dimensionless
duration shrinks in laboratory units. Their product is a bounded
action; this is not a bounded-height or fixed-duration protocol.

The basic limit and two useful parameters are
\begin{equation}
 r\downarrow0,\quad s\to\infty,\quad r\sqrt{s}\to0,
 \qquad N=\frac1{r\sqrt{s}},\quad \Xi=\frac{s\ee^{-s}}r.
                                                               \label{eq:regime}
\end{equation}
$N$ measures the growing thermal mode count. The input purity obeys
$\Tr\rho_{r,s}^2\sim\sqrt{2\pi}\,r\sqrt{s}$ and hence tends to zero.
The parameter $\Xi$ compares activated population in the raised region
with the boundary contribution to a rare output. Some results below
require $\Xi$ bounded or convergent; the resonant sharp-wall
purification result does not.

\section{One spectral measure for the boundary response}
\label{sec:source}

In the boundary coordinate $x=(y-1)/r$ the interval is
$(-\ell,\ell)$, $\ell=1/r$, and the kinetic term is $-\partial_x^2$.
For a smooth plateau profile let $f_h$ be the positive decaying
zero-energy solution on $[-d,\infty)$,
\begin{equation}
 -f_h''+hFf_h=0,\quad f_h(-d)=1,\quad m_h=-f_h'(-d)>0.
                                                               \label{eq:zeroenergy}
\end{equation}
Set $g=Ff_1$, extended by zero to $x<-d$, and
$B_\infty=-\partial_x^2+4F$ on the full line. Define
\begin{equation}
 \nu_F(\dd E)=\frac9{m_1^2}E^{-2}
       \ip{g}{\one_{\dd E}(B_\infty)g},\qquad E>0.
                                                               \label{eq:source}
\end{equation}
This is a positive spectral measure, not a probability distribution.
Its mass is infinite at zero, while its first and second moments are
finite. Write
\begin{align}
 D_F&=m_1^{-1}-m_4^{-1},\nonumber\\
 M_F&=3\dot m_1/m_1^2-D_F=\int E\,\nu_F(\dd E),\label{eq:moments}\\
 Q_F&=9\norm{g}^2/m_1^2=\int E^2\,\nu_F(\dd E),\nonumber
\end{align}
where a dot denotes differentiation with respect to wall height.
The two response functions used throughout are
\begin{equation}
 \begin{gathered}
 C_F(t)=\int\sin^2(tE/2)\,\nu_F(\dd E),\qquad\\
 \Phi_F(c)=\int\min(E,c)\,\nu_F(\dd E).          \label{eq:responses}
 \end{gathered}
\end{equation}

The local density at threshold is
\begin{equation}
 \nu_F'(E)=\frac{D_F^2}{\pi}E^{-3/2}[1+O_F(E)]
                  \quad(E\downarrow0).                       \label{eq:threshold}
\end{equation}
It follows that $C_F(t)$ is finite at every fixed time and
$0<\Phi_F(c)\le M_F$ for $c>0$, with strict upper inequality for
$0<c<4$. The coefficient $D_F$ is
unchanged by moving the reference point $-d$ farther into the free
plateau. It satisfies $D_F\ge1/2$, strictly for a genuine smooth
transition. The resulting ordered limits are
\begin{equation}
 \Phi_F(1/a)\sim\frac{4D_F^2}{\pi\sqrt a},\qquad
 C_F(t)\sim\frac{D_F^2}{\sqrt{2\pi}}\sqrt{|t|},                 \label{eq:ordered}
\end{equation}
as $a\to\infty$ or $|t|\to\infty$ after forming the continuum source.
These are not joint growing-time statements for the finite device.

For the sharp wall $m_h=\sqrt h$, $D_F=1/2$, $M_F=1$, $Q_F=9/2$,
and we omit the subscript $F$. Matching the full-line Green function
gives the explicit density
\begin{equation}
 \nu'(E)=
 \begin{cases}
 \displaystyle\frac9{4\pi E^{3/2}(1+\sqrt{4-E})^2},&0<E<4,\\[5pt]
 \displaystyle\frac{9(1+\sqrt E\sqrt{E-4})}
 {\pi E^2(\sqrt E+\sqrt{E-4})(E-3)^2},&E>4.
 \end{cases}                                                   \label{eq:sharpdensity}
\end{equation}
There is no atom at four. The apparent change there is the opening of
the right propagating channel, not an additional bound state.

\section{Coherence retained by boundary interference}
\label{sec:herald}

For a low initial eigenstate $u_E$, the part in the controlled region
is a scalar boundary value times a common analytic family:
$Fu_E=b_Eg_{E,\ell}$. Duhamel's identity therefore reads exactly
\begin{equation}
 \begin{gathered}
 K_tu_E=b_E\ee^{-itE}V_{E,\ell}(t),\qquad\\
 V_{E,\ell}(t)=\frac32 h_t(B_\ell-E)g_{E,\ell},\quad\\
 h_t(z)=\frac{\ee^{-itz}-1}{z}.                 \label{eq:column}
 \end{gathered}
\end{equation}
The value at zero is $h_t(0)=-it$. Low thermal energies approach zero,
so these columns become parallel even though their input eigenstates
are orthogonal. Their complete thermal covariance has leading term
\begin{equation}
 K_t\rho_{r,s}K_t^*\big|_{\text{low input}}
 =\frac r{s m_1^2}\ket{V_{0,\ell}(t)}\bra{V_{0,\ell}(t)}
                       +o_1(r/s),                            \label{eq:covariance}
\end{equation}
where $o_1$ is in trace norm. The source vector converges strongly to
$V_t=(3/2)h_t(B_\infty)g$, with
$\norm{V_t}^2/m_1^2=C_F(t)$.

If $\Xi\to x<\infty$, the full, off-resonant probability satisfies
\begin{equation}
 \frac{s}{r}p_{r,s}(t)\longrightarrow
       C_F(t)+x\sin^2(3t/2),                                  \label{eq:timeyield}
\end{equation}
uniformly on each fixed compact time interval. The second term is
activated volume, not another coherent source. For a fixed $t\ne0$,
put $b=x\sin^2(3t/2)$. Its positive covariance has bounded limiting
trace and vanishing operator norm after scaling by $s/r$. Consequently
\begin{equation}
 \begin{gathered}
 \lambda_1(\sigma_t)\to\frac{C_F(t)}{C_F(t)+b},\quad\\
 \lambda_2(\sigma_t)\to0,\quad\\
 \Tr\sigma_t^2\to\left[\frac{C_F(t)}{C_F(t)+b}\right]^2.
                                                               \label{eq:puritylimit}
 \end{gathered}
\end{equation}
When $b>0$, the output entropy diverges: a nonzero weight is distributed
among a growing number of small eigenvalues. When $b=0$, purity alone
does not decide entropy.

At the fundamental resonance $\tau=2\pi/3$, the activated plateau
acquires the same branch phase. For the sharp wall, cancellation leaves
only boundary terms at every initial energy, including at the barrier
threshold. In the whole regime~\eqref{eq:regime},
\begin{equation}
 p_{r,s}(\tau)\sim C(\tau)\frac r s,\qquad
 \norm{\sigma_\tau-\Pi_{r,\tau}}_1\to0,
 \quad S(\sigma_\tau)\to0,                                  \label{eq:resonance}
\end{equation}
and all fixed positive-order R\'enyi entropies tend to zero.
$\Pi_{r,\tau}$ projects onto the actual finite-box zero-energy source.
The entropy conclusion uses all-index spectral-tail bounds, given
in Appendix~\ref{app:control}, rather than a finite-dimensional
continuity inequality. The smooth plateau version has the same yield
$C_F(\tau)r/s$ and entropy conclusion, with profile-dependent constants.

The average number of fresh preparations per retained state is
$s/[C_F(\tau)r]$ to leading order. Each preparation includes a pure
ancilla and a readout; the retained purity is conditional. The initial
state becomes more mixed while the successful output becomes pure,
but the success probability tends to zero.

\subsection{Changing the controlled boundary changes the output rank}

The source picture can be tested without changing the input ensemble.
On $(0,3)$ let
\begin{equation}
 A_r=-r^2\partial_y^2+Q_L+Q_R,\quad
 Q_L=\one_{(0,1)},\quad Q_R=\one_{(2,3)}.                       \label{eq:twowell}
\end{equation}
Compare control of $3Q_R$ with control of $3(Q_L+Q_R)$, at the same
resonance and with the same complete Gibbs preparation. The limits are
\begin{center}
\begin{tabular}{lcc}
\toprule
Controlled walls & Right only & Both\\
\midrule
$p/(r/s)$ & $C(\tau)$ & $2C(\tau)$\\
Nonzero output weights & $1$ & $1/2,1/2$\\
Output purity & $1$ & $1/2$\\
Output entropy & $0$ & $\log2$\\
\bottomrule
\end{tabular}
\end{center}
Both initial parity sectors are present in each experiment. With one
controlled wall they feed the same boundary vector; with both controlled
walls they feed orthogonal final parity sectors of equal limiting
weight. This is a change of the physical controlled Hamiltonian, not
a parity postselection. The calculation in Appendix~\ref{app:switch}
keeps the factor of two in both the boundary amplitudes and the total
preparation probability.

\subsection{Repeated retention and the price of imperfect interference}

For any fixed positive integer $m$, $m$ ideal successive dark outcomes
apply $K_\tau^m$ to the same particle. The limiting source is
$K_{\infty,\tau}^{m-1}V_\tau$, and
\begin{equation}
 p_m\sim C_{m,F}\frac r s,\qquad
 C_{m,F}=\frac{\norm{K_{\infty,\tau}^{m-1}V_\tau}^2}{m_1^2}>0.
                                                               \label{eq:repeat}
\end{equation}
The normalized state again has vanishing entropy. The positivity in
this formula follows from injectivity of the actual first-resonance
Kraus operator, not a presumed dominant eigenvector of a commuting
filter. No interchange with $m\to\infty$ is needed.

Real visibility loss is described by the completely positive map
\begin{equation}
 \begin{gathered}
 \mathcal J_\nu(X)=\nu KXK^*+(1-\nu)\mathcal D(X),\qquad\\
 \mathcal D(X)=\frac{U_BXU_B^*+U_AXU_A^*}{4},                  \label{eq:visibility}
 \end{gathered}
\end{equation}
with $0\le\nu\le1$ and $\Tr\mathcal D(X)=\Tr X/2$.
For two independent stages, with $\nu_j\in[0,1]$ and
$\epsilon_j=1-\nu_j$, sufficient source-retention conditions are
\begin{equation}
 \epsilon_1+\epsilon_2\to0,\qquad
 \epsilon_1\epsilon_2=o(r/s).                                \label{eq:twonoise}
\end{equation}
For vanishing entropy the second condition strengthens to
$\epsilon_1\epsilon_2\log N=o(r/s)$. Under Eq.~\eqref{eq:twonoise},
this strengthened product condition is also necessary. The
all-false branch carries a mixture of four unitary images of the
full Gibbs state and hence an entropy $\log N+O(1)$. Independent
repetition suppresses that branch quadratically, but cannot make its
preparation or entropy cost disappear.

\section{Where dissipated work becomes population or coherence}
\label{sec:thermo}

Let $\gamma_B=\ee^{-sB_r}/Z_4$ and let $\Delta_B$ completely dephase in
the final energy basis. With natural logarithms, set
\begin{equation}
 \begin{gathered}
 I=\Srel(\rho\Vert\gamma_B),\quad\\
 P=\Srel(\Delta_B\rho\Vert\gamma_B),\quad\\
 J=\Srel(\rho\Vert\Delta_B\rho),\qquad I=P+J.                  \label{eq:allocation}
 \end{gathered}
\end{equation}
The identity is the established population--coherence decomposition
\cite{Francica2019}. Here it applies to the unselected state immediately
after a sudden change of Hamiltonian, independently of whether a Ramsey
outcome is retained. $I$ equals the mean dissipated work multiplied by
the physical inverse temperature.

Write $\Lambda=\log(1/r)$. On the paths
\begin{equation}
 \frac{s}{\Lambda}\to a\in[1,\infty),\qquad
 \Xi\to x\in[0,\infty),                                    \label{eq:thermopath}
\end{equation}
where both $a$ and $x$ are finite, the complete smooth-profile result is
\begin{equation}
 \frac{(I,P,J)}r\longrightarrow
 \bigl(M_F+3x,\ M_F-\Phi_F(1/a)+3x,\ \Phi_F(1/a)\bigr).
                                                               \label{eq:thermolaw}
\end{equation}
Necessarily $x=0$ if $a>1$; for example $s=a\Lambda$ realizes this
case. At $a=1$, any $x>0$ is realized by
$s=\Lambda+\log\Lambda-\log x$, while
$s=\Lambda+2\log\Lambda$ realizes $x=0$.
The activated sector changes the total and
population terms by $3x$, while the leading coherence coefficient is
unchanged. Its final-energy weight concentrates at four; its work
increment concentrates at three. These are different variables.

The cap $1/a$ has a direct origin. At a fixed final energy $E$, the
source population is polynomially small in the boundary parameter,
whereas the equilibrium weight is proportional to $\ee^{-sE}$.
Their competition occurs at $E=1/a$. Below it, population relative
entropy is negligible at order $r$; above it, the contribution is
$(E-1/a)\nu_F(\dd E)$. This interpretation alone would not control
the moving crossover or the infinitely many final levels. The proof
in Appendix~\ref{app:entropy} pays both with positive diagonal excess
and an energy-weighted tail estimate.

For the sharp wall, $M_F=1$ and the coherent fraction of the leading
source dissipation is $\Phi(1/a)$. At $a=1$,
\begin{equation}
             0.4040654186<\Phi(1)<0.4040654188.                \label{eq:phinumeric}
\end{equation}
For a smooth profile, $D_F>1/2$ increases the sequential large-$a$
coefficient in Eq.~\eqref{eq:ordered}. This compares that coefficient,
not all finite-$a$ curves. A useful profile-independent connection to
the retained yield is
\begin{equation}
          C_F(2\pi/3)\le a\Phi_F(1/a)\le aM_F.               \label{eq:yieldcoherence}
\end{equation}
The first inequality follows pointwise from
$\sin^2(\pi E/3)\le\min(E,1)\le a\min(E,1/a)$.

Although $I=O(r)$ tends to zero, its dimensional work need not.
Equation~\eqref{eq:units} gives
$\langle W_{\rm diss}\rangle\sim E_0(M_F+3x)/(rs)$.
The dimensionless relative entropy is not a fixed-energy cooling
efficiency. Conditional retained states and unselected work statistics
answer different operational questions.

\section{A regular work limit and a singular absolute-work law}
\label{sec:work}

Let $(E_n,u_n)$ and $(\lambda_j,\phi_j)$ be the complete initial and
final eigenpairs of Eq.~\eqref{eq:ham}. The sudden-quench work law in
barrier units is
\begin{equation}
 \Pi_{r,s}=\frac1{Z_1}\sum_{n,j}\ee^{-sE_n}
        |\ip{\phi_j}{u_n}|^2\delta_{\lambda_j-E_n}.            \label{eq:tpm}
\end{equation}
It is the two-projective-measurement work distribution. In our retained
state protocol those energy measurements are not made: the spectral
formula is also an exact expansion of the two controlled propagators.
In particular,
\begin{equation}
       p_{r,s}(t)=\int\sin^2(t\omega/2)\,\Pi_{r,s}(\dd\omega).
                                                               \label{eq:ramseytpm}
\end{equation}

Under Eq.~\eqref{eq:regime} with $\Xi\to x<\infty$, the finite positive
measures $\zeta_{r,s}(\dd\omega)=(s/r)\omega^2\Pi_{r,s}(\dd\omega)$
converge against every bounded continuous test function:
\begin{equation}
 \zeta_{r,s}\ \Longrightarrow\ 
       \omega^2\nu_F(\dd\omega)+9x\delta_3.                  \label{eq:weightedlaw}
\end{equation}
The first term is supported on positive work. The limiting second
moment is $Q_F+9x$. Dividing the sine-square test function in
Eq.~\eqref{eq:ramseytpm} by $\omega^2$, with its continuous value at
zero, proves Eq.~\eqref{eq:timeyield}. Absolute work, however, requires
the singular quotient $1/|\omega|$. Equation~\eqref{eq:weightedlaw}
cannot be used for that quotient.

For the sharp wall, direct finite-spectrum analysis resolves it.
For every fixed $c>0$, on Eq.~\eqref{eq:regime} with $\Xi=O(1)$,
\begin{align}
 \frac s r\mathbb E\min(|\omega|,c)
 &=\Phi(c)+\frac12+\frac{\log N}{\pi^{3/2}\sqrt s}
                   \nonumber\\
 &\quad+\min(3,c)\Xi+o(1),\label{eq:abscap}\\
 \frac s r\mathbb E|\omega|
 &=\frac32+\frac{\log N}{\pi^{3/2}\sqrt s}+3\Xi+o(1),
                                                            \label{eq:absolute}\\
 \frac s r\mathbb E\omega&=\frac32+3\Xi+o(1),\qquad\nonumber\\
 \frac s r\mathbb E(-\omega)_+
       &=\frac{\log N}{2\pi^{3/2}\sqrt s}+o(1).                \label{eq:negative}
\end{align}
The remainders are additive even when the displayed logarithmic term
diverges. The $1/2$ term comes from paired low levels whose work
increments shrink to zero. The logarithm comes from a harmonic sum
of transitions between nearby but differently indexed low levels.
Neither term is obtained by applying a singular test function to
Eq.~\eqref{eq:weightedlaw}. Appendix~\ref{app:infrared} derives them
before taking the continuum limit.

Physical work is $W=(E_0/r^2)\omega$. The following consequence is
independent of the unit used for energy. If
$s/\Lambda\to a>1$, then $\Xi\to0$ and
\begin{equation}
 \frac{\mathbb E|W|}{\mathbb EW}
       \sim\frac{2\sqrt{s}}{3a\pi^{3/2}},\qquad
 \frac{\mathbb E(-W)_+}{\mathbb E(W)_+}\longrightarrow1.       \label{eq:ratio}
\end{equation}
Large positive and negative absolute components increasingly cancel
in the positive net work. Negative individual two-measurement outcomes
are compatible with the nonnegative potential increment $B-A=3F$;
the mean is $3\Tr\rho F\ge0$. There is no violation of passivity or
of the positivity of dissipated work.

The near-elastic sector is infrared \emph{relative to the barrier
energy}. It need not have small laboratory energy. For example, a
typical adjacent spacing is of order $r/\sqrt s$ in $\omega$, hence
$E_0/(r\sqrt s)$ in $W$, which grows in Eq.~\eqref{eq:regime} at
fixed $E_0$. Similarly, fixed $t$ in Eq.~\eqref{eq:timeyield} means a
shrinking laboratory duration. Equations~\eqref{eq:absolute}--
\eqref{eq:ratio} avoid interpreting a choice of units as a new
low-frequency laboratory effect.

\section{Recovering a coherent coefficient from actual counts}
\label{sec:counts}

For any positive measure on $(0,\infty)$ with finite first moment,
cosine integration gives the positive reconstruction
\begin{equation}
 \Phi_F(c)=\frac4\pi\int_0^\infty
              \frac{1-\cos(ct)}{t^2}C_F(t)\,\dd t.            \label{eq:positive}
\end{equation}
Its kernel has total mass $2c$. Thus an all-time uniform additive
error $\eta$ in $C_F$ changes the reconstructed coefficient by at
most $2c\eta$. This stability concerns a fixed capped moment;
identifying the whole measure by differentiating cap data is a
different inverse problem. Finite-resolution work measurements and
coherence inference from energy data have their own established
frameworks \cite{Ahmad2022,Alonso2026}. Here the target is the leading
known-model continuum coefficient, not exact finite-state coherence
tomography or a replacement for measurement backaction.

The finite work law has the same integral identity with
$C_F(t)$ replaced by $p_{r,s}(t)$ and $\Phi_F(c)$ by
$\mathbb E\min(|\omega|,c)$. Equations~\eqref{eq:abscap} and
\eqref{eq:timeyield} show why the scaled finite integral and the
continuum-source integral cannot be interchanged. An infrared-rejected
functional gives a regular bridge: for fixed $0<\eta<c$,
\begin{align}
 \frac s r\mathbb E\{\min(|\omega|,c)-\min(|\omega|,\eta)\}
 \nonumber\\
 \longrightarrow\Phi_F(c)-\Phi_F(\eta)
             +x[\min(3,c)-\min(3,\eta)].                    \label{eq:reject}
\end{align}
Its quotient by $\omega^2$ is bounded and continuous. The device limit
is taken before $\eta\downarrow0$. The corresponding time kernel is
$4[\cos(\eta t)-\cos(ct)]/(\pi t^2)$ and need not be positive.

For a finite positive source quadrature, let $t_i$ be the midpoints
of a partition of $[0,T]$ and
\begin{equation}
 \begin{gathered}
 w_i=\frac4\pi\int_{I_i}\frac{1-\cos(ct)}{t^2}\,\dd t,
 \qquad\\ W=\sum_iw_i\le2c,\qquad \lambda=s/r.                 \label{eq:weights}
 \end{gathered}
\end{equation}
The actual device quantity is $S=\lambda\sum_iw_ip_{r,s}(t_i)$.
Choose node $i$ with probability $w_i/W$, independently on each
fresh preparation, and record whether the dark outcome occurs.
For independent trials with the same node response probabilities,
after $n$ trials the pooled count is exactly binomial with parameter
$\bar p=S/(\lambda W)$. The estimator is
$\widehat S=\lambda W k/n$. If a deterministic full-device
certificate gives $d_-\le S-\Phi_F(c)\le d_+$, any exact binomial
confidence interval $[L_n(k),U_n(k)]$ gives
\begin{equation}
 \Phi_F(c)\in[\lambda W L_n(k)-d_+,\,
              \lambda W U_n(k)-d_-].                         \label{eq:confidence}
\end{equation}
There is no subtraction of a fitted model mean from the raw counts.

A convenient sufficient preparation budget, when $S\le D$ is
certified, is
\begin{equation}
 n\ge\frac{2\lambda W\log(2/\alpha)}{\epsilon^2}
                      (D+\epsilon/3).                        \label{eq:shots}
\end{equation}
It guarantees $|\widehat S-S|\le\epsilon$ with probability at least
$1-\alpha$. The rare-event cost is proportional to $s/r$, not its
square. It nevertheless remains large. For the sharp wall,
$r=1/200$, $s=128$, $c=1$, $T=32$, and 256 midpoint nodes of spacing
$1/8$, certified continuum enclosures give two distinct designs:
\begin{center}
\begin{tabular}{lrr}
\toprule
Systematic accounting & Preparations & Total error\\
\midrule
Signed, fixed device & $62\,217\,817$ & $<0.050480$\\
Separate absolute errors & $5\,999\,715\,692$ & $<0.099933$\\
\bottomrule
\end{tabular}
\end{center}
Both have confidence at least $95\%$. The second is less than
$\Phi(1)/4$. The first uses opposing source-grid and finite-device
biases at the certified point; the second pays their absolute sizes
separately. These counts include unsuccessful trials, fresh Gibbs
states, pure ancillas, and readouts. They are sufficient mathematical
budgets, not observed data, optimal lower bounds, or a claim of
practical experimental advantage.

Resonance-only sampling does not generally determine the capped
coefficient: $\sin^2(n\tau(E+3)/2)=\sin^2(n\tau E/2)$. For example,
adding respectively $\beta\delta_{7/2}$ and
$\beta(\delta_{1/2}+\delta_{13/2})/2$ to a common positive source
preserves the first moment and every resonant sample but changes
$\Phi(1)$ by $\beta/4$. This is an obstruction in the positive-measure
class, not a claim to have constructed two actual wall profiles with
identical resonant data. Off-resonant measurements and their activated
background must be treated explicitly.

\section{A finite smooth device with unrecorded timing error}
\label{sec:finite}

The asymptotic results do not specify an experimental onset. A
separate validated calculation includes the actual analytic wall
\begin{equation}
 \begin{gathered}
 F_r(y)=\tfrac12\{1+\operatorname{erf}[(y-1)/0.025]\},\quad\\
 r=0.05,\quad s=8,\quad\\
 A=-r^2\partial_y^2+F_r,\quad B=A+3F_r.
                                                               \label{eq:erf}
 \end{gathered}
\end{equation}
Its 10--90 percent transition width is $0.9061938024\ldots$ times
the penetration length $rL_0$; it is not a vanishingly small
perturbation of a sharp wall. The accepted output is averaged over
any input-independent, unrecorded duration distribution supported
on $2\pi/3\pm0.001$. Visibility is $\nu=0.9999$ in
Eq.~\eqref{eq:visibility}. Readout independently has false-positive
probability $10^{-5}$ conditional on a bright outcome, and
false-negative probability $0.01$ conditional on a dark outcome.
For precisely this model, outward-rounded enclosures give
\begin{align}
 \Tr\rho^2&<0.386060,\nonumber\\
 \Tr\sigma_{\rm recorded}^2&>0.9139989,\label{eq:certificate}\\
 0.00257508&<p_{\rm recorded}<0.00261934.\nonumber
\end{align}
The mean number of fresh preparations per recorded output is below
389. This expectation is not a high-probability guarantee or a
tomography cost. For independent attempts, a lower bound $p_-$
gives at least one success with probability $1-\alpha$ after
$n\ge\log\alpha/\log(1-p_-)$ attempts.

The timing estimate is in the laboratory Schr\"odinger frame.
For the actual Gibbs moments,
\begin{equation}
 \begin{gathered}
 \norm{\partial_t K_t\sqrt\rho}_2
 \le\sqrt{\Tr\rho A^2}+\frac32\sqrt{\Tr\rho F_r^2}
 \\=L<0.220728.                                                \label{eq:jitter}
 \end{gathered}
\end{equation}
A purification retaining a classical duration register bounds the
unnormalized amplitude error by $L\Delta t$; normalization by the
rare success probability is then paid explicitly. A fixed-time
interaction-picture bound alone would not establish this statement
about a mixture of unrecorded durations.

The certificate uses validated ODE integration for the smooth
transition, certified eigenvalue labels, complete Gibbs tails, and
complete final spectral tails. Computational prefixes are not
physical energy cutoffs. Appendix~\ref{app:certificate} gives the
enclosure construction and its noise propagation; the public
verification scripts reproduce the reported inequalities.

Physical resolution, temperature, and duration cannot be chosen
independently. If the 10--90 edge width is $W=a_{\rm erf}rL_0$,
where $a_{\rm erf}=0.9061938024\ldots$, then
\begin{equation}
 \begin{gathered}
 V=\frac{\hbar^2a_{\rm erf}^2}{2mW^2},\quad
 k_BT=V/s,\quad\\ t_{\rm phys}=\tau\hbar/V,\quad
 k_BTt_{\rm phys}/\hbar=\tau/s.                              \label{eq:tradeoff}
 \end{gathered}
\end{equation}
For $^{87}$Rb and a $400$-nm width, this example corresponds to
approximately $1.79$ nK and $1.12$ ms; the timing half-width is
about $0.533\,\mu$s. This conversion fixes a demanding target,
not an apparatus demonstration. Finite outer walls, transverse
excitations, state-dependent losses, preparation and reset overhead,
and interactions require their own experimental error budgets.

\section{Discussion}

The common source measure organizes three physically different
questions. A coherent boundary experiment retains a nearly pure
vector because low-energy columns align. Unselected thermodynamics
splits its leading source dissipation at a Gibbs crossover, producing
a capped moment. A quadratic work weighting has a regular continuum
limit, whereas the less regular absolute-value functional retains
the discrete nearby-level sector. These statements are compatible
precisely because they do not exchange conditioning, thermal limits,
and singular spectral testing.

The finite smooth example makes conditional purification a quantitative
prediction for a specified channel. The count construction addresses
a different question: how many complete preparations suffice to infer
a capped source coefficient. Its large cost is part of the result.
Neither larger conditional purity nor a finite asymptotic coefficient
alone establishes a measurement advantage.

Several experimentally meaningful comparisons remain possible within
this framework: changing which wall is controlled at fixed preparation,
measuring the dependence on visibility and timing spread, and comparing
regular quadratic work statistics with the near-elastic contribution
to absolute work. For smooth profiles the source and thermodynamic
laws are proved directly; the sharp-wall logarithmic absolute-work
coefficient is a separate exact matching result. Extending that
coefficient to a prescribed smooth family would require a uniform
two-energy scattering expansion, not merely its zero-energy source
density. This distinction keeps the present predictions tied to the
Hamiltonians for which they have been established.

\begin{acknowledgments}
The author declares the use of consumer-grade Gemini-3 and GPT-6 models for literature search, exploring and developing ideas, codes and proofs, drafting and LaTeX typesetting; the author takes full and sole responsibility for all contents of this work.
\end{acknowledgments}

\section*{Data availability}
Numerical results reported here are reproducible from the verification
scripts and certificate data in the repository - Repository - 
\href{https://github.com/arun-chandru/coherent_work}{github.com/arun-chandru/coherent\_work}.

\appendix

\section{Boundary source and regular work limit}
\label{app:source}

We give the estimates used by the thermodynamic and work results,
including the complete thermal sectors. Throughout this appendix
$\ell=1/r$, $b=\ell-d$, and $\mathcal L=\log(1/r)$; the latter must
be distinguished from the boundary-coordinate length. Constants may
depend on the fixed profile and on fixed energy cutoffs. For the sharp
wall take $d=0$ and use its elementary sine and hyperbolic-sine solutions.
Write $A_{h,\ell}=-\partial_x^2+hF$ with Dirichlet conditions on
$(-\ell,\ell)$, and abbreviate $A_\ell=A_{1,\ell}$ and
$B_\ell=A_{4,\ell}$.

\subsection{Finite-box normalization and boundary amplitudes}

Dirichlet and Neumann splitting at $-d,d$ gives, uniformly for
$1\le h\le4$,
\begin{equation}
 N_{h,\ell}(E)=\frac b\pi\bigl(\sqrt E+\sqrt{(E-h)_+}\bigr)
                      +O_d(1+\sqrt E).                 \label{eq:apcount}
\end{equation}
The error is bounded by the free count on the fixed middle interval
and the endpoint integer errors. The same bound holds for strict
counts. Stieltjes integration therefore yields
\begin{align}
 Z_h&=\frac{b(1+\ee^{-hs})}{2\sqrt\pi\sqrt s}
                              +O_d(1+s^{-1/2}),\nonumber\\
 T_h&:=\Tr\one_{[h,\infty)}(A_{h,\ell})\ee^{-sA_{h,\ell}}
 \nonumber\\
 &=\frac{b\ee^{-hs}}{2\sqrt\pi\sqrt s}
       [1+O_d(s^{-1/2}+\sqrt s/b)].                    \label{eq:apheat}
\end{align}
For the second expression integrate the increment from $h-$: the
right square root gives the displayed term, the left increment costs
$O(b\ee^{-hs}/s)$, and the counting error costs $O_d(\ee^{-hs})$.
Thus possible threshold eigenvalues are included. The same argument
gives, for every fixed $e_0>0$,
\begin{equation}
 Z_1^{-1}\Tr\one_{[e_0,\infty)}(A_\ell)\ee^{-sA_\ell}
                      \le C_{F,e_0}\ee^{-se_0}.         \label{eq:aptail}
\end{equation}
In the basic regime, $Z_h\sim(2\sqrt\pi r\sqrt s)^{-1}$; applying
this at $s$ and $2s$ proves the input-purity formula in the main text.

The right Dirichlet operator $H_{h,D}$ on $(-d,\ell)$ has a uniform
gap $\gamma_d=1/\max(1,2d^2)$: for $u(-d)=0$,
$\int_{-d}^d|u|^2\le2d^2\int_{-d}^d|u'|^2$, while $hF\ge1$
on $(d,\ell)$. Let $f_{h,E,\ell}$ equal one at $-d$, vanish at
$\ell$, and solve the right equation; set
$m_{h,\ell}(E)=-f'_{h,E,\ell}(-d)$. Below a fixed small energy,
\begin{equation}
 \begin{aligned}
 f_{h,E,\ell}&=(I-EH_{h,D}^{-1})^{-1}f_{h,0,\ell},\qquad\\
 \partial_E m_{h,\ell}(E)&=-\norm{f_{h,E,\ell}}^2.
 \end{aligned}                                         \label{eq:apode}
\end{equation}
The second identity follows by differentiating the equation and
integrating its Wronskian. The zero-energy Dirichlet principle bounds
$m_{h,\ell}(0)\ge(2d+1)^{-1}$: discard the transition potential and
minimize $|1-u(d)|^2/(2d)+|u(d)|^2$. Zero-energy solutions have bounded
norm, and the terminal constant plateau makes them and their
low-energy derivatives converge exponentially as $\ell\to\infty$.
Matching to the left sine, then integrating the normalization, gives
\begin{align}
 kb+\arctan\frac{k}{m_{h,\ell}(k^2)}&=n\pi,\nonumber\\
 |u_E(-d)|^2
 &=\frac{2E}{\begin{gathered}
 b[E+m_{h,\ell}(E)^2]+m_{h,\ell}(E)\\
                       -2E\partial_E m_{h,\ell}(E)
 \end{gathered}}\nonumber\\
 &=\frac{2E}{b m_h^2}
       [1+O_F(E+b^{-1}+\ee^{-c_Fb})].                    \label{eq:apamp}
\end{align}
The phase count has bounded discrepancy from $bk/\pi$. For fixed
nonnegative integers $j$, Gaussian Stieltjes integration implies
\begin{equation}
 \sum_{E_n<\delta}\ee^{-sE_n}|u_n(-d)|^2E_n^j
 =\frac{\Gamma(j+3/2)}{\pi m_h^2}s^{-j-3/2}[1+o(1)],    \label{eq:apmom}
\end{equation}
where $\delta>0$ is fixed and sufficiently small. Relative errors are
bounded by a constant times
$s^{-1}+\sqrt s/b+\ee^{-\delta s/2}+\ee^{-c_Fb}$, uniformly in $h$.

We shall also use the all-mode subthreshold estimates
\begin{equation}
 \begin{aligned}
 \sum_{E_n<h}\norm{F^{1/2}u_n}^2&\le C_F,\qquad\\
 \sum_{E_n<h}\ee^{-sE_n}\norm{Fu_n}^2&=O_F(s^{-3/2}).
 \end{aligned}                                          \label{eq:apcollar}
\end{equation}
For low energies these follow from Eq.~\eqref{eq:apamp}. For
$\delta\le E<h$, left sine normalization is $O(b^{-1/2})$, and a
bounded transfer matrix across the fixed transition gives the same
bound at $d$. The right-plateau mass is at most
$C_F\min\{1,[b\sqrt{h-E}]^{-1}\}$. Integrate this increasing bounded
weight against $N(E)=b\sqrt E/\pi+O_F(1)$. The principal integral is
bounded by $\int_0^h[E(h-E)]^{-1/2}\dd E=\pi$; the discrepancy is
bounded by the total variation and endpoint value. The middle interval
costs a constant by $\sum_{E_n\le4}|u_n(x)|^2\le C_F$. To prove this
local bound apply the one-dimensional evaluation inequality to
$(A_{h,\ell}+1)^{-1}$ and use its spectral expansion. This argument
includes modes arbitrarily close to $h$. Multiplying the low sector
by the heat weight and using Eq.~\eqref{eq:apmom}, then charging the
rest exponentially, proves the second assertion, also with $F$ in
place of $F^2$.

\subsection{Activated mass away from the controlled plateau}

When $b/s\to\infty$, the following local estimate suffices:
\begin{align}
 U_h&:=\Tr(1-F)\one_{[h,\infty)}(A_{h,\ell})\ee^{-sA_{h,\ell}}
 \nonumber\\
 &\le C_F\ee^{-hs}(b/s+1)+C_Fb\ee^{-(h+1)s}/\sqrt s,\nonumber\\
 U_h/T_h&=O_F(s^{-1/2}+\sqrt s/b+\ee^{-s})=o(1).         \label{eq:aplocal}
\end{align}
Here is a local resolvent proof which avoids eigenfunction localization.
On the free left plateau put $t=x+\ell\in(0,b)$, $z=k^2$, $k=p+iq$,
$p,q>0$. The right solution, normalized at $-d$, has logarithmic
derivative $m(z)=\mathcal E-z\mathcal N$, where $\mathcal E,\mathcal N$
are its nonnegative energy and norm integrals. Writing $m=u-iv$ gives
$qu+pv=q(\mathcal E+|k|^2\mathcal N)\ge0$, hence
$R(z)=[m(z)+ik]/[ik-m(z)]$ obeys $|R(z)|\le1$. The diagonal Green
function is exactly
\begin{equation}
 G_z(t,t)=\frac{i}{2k}
 \frac{1-\ee^{2ikt}+R\ee^{2ik(b-t)}-R\ee^{2ikb}}
      {1+R\ee^{2ikb}}.                                  \label{eq:apgreen}
\end{equation}
For real parts in a fixed interval above zero and imaginary part
$1/s$, one has $q\asymp1/s$. The denominator stays away from zero
when $b/s\to\infty$, so $|G_z(t,t)|\le C$. The spectral inequality
$\one_{[v-\eta,v+\eta]}(A)\le2\eta\operatorname{Im}(A-v-i\eta)^{-1}$
then gives a left-local count $Cb\eta$. Sum cells of width $1/s$ in
$[h,h+1]$ with their heat weights to obtain $Cb\ee^{-hs}/s$.
The fixed middle interval costs $C_F\ee^{-hs}$ by the preceding
local spectral bound; Eq.~\eqref{eq:aptail} pays every energy above
$h+1$. This proves Eq.~\eqref{eq:aplocal}.

For the sharp wall a frequency-bin bound removes the condition
$b/s\to\infty$. On the physical left interval every mode above one
is $a_n\sin(\omega_n y)$, $\omega_n=\sqrt{E_n}/r$. For sufficiently
small $r$, the sine norm is at least $2/3$, and its frequency derivative
has norm at most $2/3$. The normalized sine at the start of a frequency
bin of width $1/4$ thus overlaps every sine in that bin by at least
$1/2$. Bessel's inequality in the original full eigenbasis gives
$\sum_{n\text{ in bin}}|a_n|^2\le4$. Counting bins yields
\begin{align}
 \sum_{1\le E_n\le E}\norm{(1-F)u_n}^2&\le4+8(E-1)/r,\nonumber\\
 T_-:=\sum_{E_n\ge1}\ee^{-sE_n}\norm{(1-F)u_n}^2
                    &\le\ee^{-s}(4+8/(rs)).              \label{eq:apsharpactivated}
\end{align}
The last step integrates the positive local spectral count, including
a possible eigenvalue at one.

\subsection{Source moments and threshold law}

Zero-energy Wronskians give
\begin{equation}
 \dot m_h=\int Ff_h^2\dd x,\qquad
 3\int Ff_1f_4\dd x=m_4-m_1.                            \label{eq:apwronskian}
\end{equation}
Let $w=[f_1-(m_1/m_4)f_4]/3$ on the right and extend it constantly
to the left. Its derivative matches zero and $B_\infty w=g$
distributionally. Although the constant left tail generally prevents
$w\in L^2$, its homogeneous energy is finite. Slow left cutoffs add
arbitrarily little energy, while
$|\int Ff_1v|^2\le(\int Ff_1^2)(\int F|v|^2)$ makes the source
functional continuous in the energy norm. The form inverse, equivalently
the monotone negative-resolvent limit, gives, with
$\mu_g(\dd E)=\ip{g}{\one_{\dd E}(B_\infty)g}$,
\begin{equation}
 \int E^{-1}\mu_g(\dd E)
 =\ip{g}{w}
 =\frac{\dot m_1}{3}-\frac{m_1}{3m_4}\int Ff_1f_4\dd x. \label{eq:apform}
\end{equation}
Multiplication by $9/m_1^2$ proves the first-moment identity
in Eq.~\eqref{eq:moments}; the second is immediate from the source
definition. Positivity is strict because $g\ne0$.

The full-line operator has no eigenvalues: nonnegativity excludes
negative ones and the free left equation excludes an $L^2$ solution
at every nonnegative energy. For $0<E<4$, let $h_E$ be the decaying
right solution with $h_E(d)=1$, $h_E'(d)=-\sqrt{4-E}$. Gluing to the
outgoing free left solution gives
\begin{equation}
 \begin{aligned}
 \mu_g'(E)&=\frac{\sqrt E}{\pi}
 \frac{|\int_{-d}^{\infty}g h_E\dd x|^2}
      {h_E'(-d)^2+E h_E(-d)^2},\qquad\\
 \mu_g(\dd E)&=\ip{g}{\one_{\dd E}(B_\infty)g}.
 \end{aligned}                                          \label{eq:apdensity}
\end{equation}
One first derives this using a boundary-normalized right solution;
the displayed form removes apparent poles at zeros of $h_E(-d)$.
The denominator never vanishes. The numerator is real analytic and
nonzero at zero by Eq.~\eqref{eq:apwronskian}, so each nonempty
subinterval of $(0,4)$ has positive source measure. Expansion at zero
proves Eq.~\eqref{eq:threshold}.

For the sharp wall put $f=\ee^{-x}\one_{x>0}$,
$\alpha=\sqrt{-z}$, $k=\sqrt{4-z}$ with positive real parts. Direct
solution and derivative matching of $(B_\infty-z)v=f$ gives
\begin{equation}
 \ip{f}{(B_\infty-z)^{-1}f}
          =\frac{k+2+\alpha}{2(k+1)^2(k+\alpha)}.         \label{eq:apsharpres}
\end{equation}
Upper-half-plane boundary values and Stieltjes inversion give
Eq.~\eqref{eq:sharpdensity}, including the continuous value at four
and the absence of an atom there. The value at $z=0-$ is $1/9$,
independently proving $\int E\nu(\dd E)=1$. Compact-interval
inversion followed by exhaustion includes the entire high-energy tail.

Reference-point invariance follows because moving $-d$ left by $a_0$
sends $m_h^{-1}$ to $m_h^{-1}+a_0$ and leaves both $g/m_1$ and
$\dot m_h/m_h^2$ unchanged. To prove $D_F\ge1/2$, put
$u=f_h/m_h$, $v=-u'$. Then $v(-d)=1$, $v\ge0$, $v'=-hFu$, and
$(v^2-hu^2)'=2h(1-F)uv\ge0$, with terminal value zero. Thus
$u\ge v/\sqrt h$ and
\begin{equation}
 \frac{\dot m_h}{m_h^2}
   =-\frac1h\int uv'\dd x
   \ge-\frac1{h^{3/2}}\int vv'\dd x
   =\frac1{2h^{3/2}}.                                    \label{eq:apscatteringbound}
\end{equation}
Integration from one to four proves the bound, strictly when a smooth
transition has $0<F<1$ on an interval. Splitting the cap at $E=1/a$
proves the first ordered limit in Eq.~\eqref{eq:ordered}. For the
second subtract $(D_F^2/\pi)E^{-3/2}\dd E$, leaving a finite signed
measure, and use
$\int_0^\infty\sin^2(u/2)u^{-3/2}\dd u=\sqrt{\pi/2}$.

\subsection{The regular measure and fixed-time state limits}

For a bounded continuous test $f$, the exact identity
$(\lambda_j-E_n)\ip{\phi_j}{u_n}=3\ip{\phi_j}{Fu_n}$ gives
\begin{equation}
 \int f\dd\zeta_{r,s}
 =9\frac sr\sum_n\rho_n
                   \ip{Fu_n}{f(B_\ell-E_n)Fu_n}.         \label{eq:apworkidentity}
\end{equation}
At a zero gap the numerator vanishes as well. In the low sector write
$Fu_n=b_ng_{E_n,\ell}$, $b_n=u_n(-d)$. Equation~\eqref{eq:apode}
gives $\norm{g_{E,\ell}-g_{0,\ell}}\le C_FE$, and
$g_{0,\ell}\to g$ after zero extension. Increasing Dirichlet intervals
converge in strong resolvent sense. The associated vector spectral
measures are tight because their masses converge and their compactly
supported continuous tests converge. On a compact interval $f$ is
uniformly continuous, so the energy shift also converges, yielding
\begin{equation}
 \begin{gathered}
 \ip{g_{E,\ell}}{f(B_\ell-E)g_{E,\ell}}
       \longrightarrow\ip{g}{f(B_\infty)g}\\
                       \quad(E\to0,\ \ell\to\infty).
 \end{gathered}                                         \label{eq:apshift}
\end{equation}
This argument does not assume global uniform continuity of $f$.
Equations~\eqref{eq:apheat} and \eqref{eq:apmom} give
$(s/r)\sum_{E_n<\delta}\rho_n|b_n|^2\to m_1^{-2}$, with vanishing
weighted mass above any fixed positive low energy. Thus this sector
in Eq.~\eqref{eq:apworkidentity} converges to
$(9/m_1^2)\ip{g}{f(B_\infty)g}$. The rest of the subthreshold
sector is at most $C_F\norm f_\infty s^{3/2}\ee^{-s\delta}$ by
Eq.~\eqref{eq:apcollar}.

The complete activated sector is $O(\Xi)$ and hence vanishes if
$x=0$. If $x>0$, then $\ell\sim x\ee^s/s$, so
Eq.~\eqref{eq:aplocal} applies. The exact shifted identity
\begin{equation}
       (B_\ell-E_n-3)u_n=-3(1-F)u_n                       \label{eq:apshiftthree}
\end{equation}
implies concentration of its averaged work measure at three: the
fraction outside $|\omega-3|<\eta$ is bounded by $9\eta^{-2}$ times
the averaged $\norm{(1-F)u_n}^2=o(1)$. Replacing $Fu_n$ by $u_n$
in Eq.~\eqref{eq:apworkidentity} costs at most
$2\norm f_\infty\norm{(1-F)u_n}$ and vanishes by Cauchy--Schwarz.
Since $(s/r)T_1/Z_1\to x$, the result is $9xf(3)$. This proves
Eq.~\eqref{eq:weightedlaw}, including its total mass. The same low
calculation with $F$ instead of $F^2$ gives
\begin{equation}
 \frac sr\mathbb E\omega\to\frac{3\dot m_1}{m_1^2}+3x
                              =M_F+D_F+3x.               \label{eq:apsigned}
\end{equation}

The continuous test $\sin^2(t\omega/2)/\omega^2$, valued $t^2/4$
at zero, proves Eq.~\eqref{eq:timeyield}. Uniformity on fixed bounded
time intervals follows from
$(s/r)|p'_{r,s}(t)|\le(|t|/2)\zeta_{r,s}(\mathbb R)$ and a finite
time net. For the state version use Eq.~\eqref{eq:column}: the same
vector convergence and low moment give trace-norm convergence of
the scaled low covariance to $m_1^{-2}\ket{V_t}\bra{V_t}$. The
collar vanishes by $\norm{K_tu_n}\le(3/2)|t|\norm{Fu_n}$.
The scaled activated covariance has trace tending to
$b=x\sin^2(3t/2)$ and operator norm at most $\Xi/Z_1\to0$, since
$K_t$ is a contraction. Its squared Hilbert--Schmidt norm and its
cross term with the source vanish. This proves Eq.~\eqref{eq:puritylimit}.
The source norm is positive for $t\ne0$ by Eq.~\eqref{eq:apdensity}.
When $b>0$,
$S(\sigma_t)\ge-(1-\lambda_1(\sigma_t))\log\lambda_2(\sigma_t)
\to\infty$. No entropy continuity is inferred when $b=0$.

\section{Complete static population and coherence allocation}
\label{app:entropy}

We prove Eq.~\eqref{eq:thermolaw} for finite fixed $a\ge1$ and
finite $x\ge0$ in Eq.~\eqref{eq:thermopath}. These paths have
$b/s\to\infty$. All estimates below concern the full unselected
sudden-quench state.

\subsection{Total entropy and positive diagonal excess}

The low thermal moments, Eq.~\eqref{eq:apwronskian}, and activated
local estimate give, uniformly for $1\le h\le4$,
\begin{equation}
 q_h:=Z_h^{-1}\Tr F\ee^{-sA_{h,\ell}}
 =\frac r s\frac{\dot m_h}{m_h^2}[1+o(1)]
                         +\ee^{-hs}[1+o(1)].             \label{eq:apheight}
\end{equation}
Each error is relative to its own positive term. Exact Duhamel trace
differentiation, with $\kappa=Z_4/Z_1$, then gives
\begin{align}
 \log\kappa&=-s\int_1^4q_h\dd h=-rD_F+o(r),\nonumber\\
 1-\kappa&=rD_F+o(r),\qquad\nonumber\\
 I&=3sq_1+\log\kappa=rM_F+3s\ee^{-s}+o(r).                \label{eq:aptotalentropy}
\end{align}
The integrated activated term is $O(\ee^{-s})=o(r)$. Entropies are
finite at finite parameters: $\Tr B\rho<\infty$ follows from the
bounded perturbation and the initial Gibbs mean energy; a Gibbs
variational bound gives finite final diagonal entropy. Thus $I=P+J$
is an equality between finite quantities.

Write $p_j=\ip{\phi_j}{\rho\phi_j}$,
$q_j=Z_4^{-1}\ee^{-s\lambda_j}$, $t_j=p_j/q_j$, and
$d_j=p_j-\kappa q_j$. Scalar Jensen in the spectral measure of $A$
for $\phi_j$ gives
\begin{equation}
 \begin{gathered}
 d_j\ge0,\quad t_j\ge\kappa,\quad
 \sum_jd_j=1-\kappa=O(r),\quad\\
 \sum_j|p_j-q_j|\le2(1-\kappa)=O(r).
 \end{gathered}                                         \label{eq:apexcess}
\end{equation}
Indeed $\langle A-\lambda_j\rangle_j=-3\langle F\rangle_j$ and
$\langle(A-\lambda_j)^2\rangle_j=9\langle F^2\rangle_j$ exactly.
No operator-monotonicity assertion for exponentials is used. Put
$G(t)=t\log t-t+1$. Then $P=\sum_jq_jG(t_j)$ has nonnegative
summands and $G(t)\le(t-1)\log t$ for $t\ge1$. The whole sector
$\kappa\le t_j<1$ costs $O((1-\kappa)^2)=O(r^2)$.

\subsection{Low final energies and the moving crossover}

For every fixed $c<4$,
$\sup_{\lambda_j\le c}\langle F\rangle_j\le C_{F,c}r$.
At small energy this is Eq.~\eqref{eq:apamp}; above a small cutoff
left sine normalization, compact transfer, and uniform right-plateau
decay prove it. Taylor's scalar exponential bound and the two exact
moments therefore imply, for $c<1/a$,
\begin{equation}
 \kappa\le t_j\le\kappa[1+C_{F,c}r(s+s^2\ee^{sc})]=1+o(1)
                          \quad(\lambda_j\le c).          \label{eq:aplowentropy}
\end{equation}
Since $G(t)=o(1)|t-1|$ uniformly on this neighborhood of one,
Eq.~\eqref{eq:apexcess} makes the complete low sea contribute $o(r)$.

Put $t_0=1/a$ and $\mathcal K=[3t_0/4,5t_0/4]\subset(0,4)$.
We require the pointwise bound
\begin{equation}
             0\le d_j\le C_{F,a}r^2/s
                              \quad(\lambda_j\in\mathcal K). \label{eq:apcollarexcess}
\end{equation}
The nonnegative scalar remainder
$R_{s,\lambda}(E)=\ee^{-sE}-\ee^{-s\lambda}
+s\ee^{-s\lambda}(E-\lambda)$ gives the exact identity
\begin{align}
 \langle\ee^{-sA}\rangle_j-\ee^{-s\lambda_j}
 &=3s\ee^{-s\lambda_j}\langle F\rangle_j\nonumber\\
 &\quad+\sum_nR_{s,\lambda_j}(E_n)|\ip{\phi_j}{u_n}|^2. \label{eq:apremainder}
\end{align}
Choose a fixed low cutoff $\delta$ below $\mathcal K$. For $E<\delta$,
$R\le\ee^{-sE}$ and
\begin{equation}
 (\lambda_j-E)\ip{\phi_j}{u_E}
       =3u_E(-d)\ip{\phi_j}{Ff_{1,E,\ell}}.              \label{eq:apbandidentity}
\end{equation}
The denominator is separated from zero; compact transfer and plateau
decay bound the last inner product squared by $Cr$. Equation
\eqref{eq:apmom} therefore pays this entire low sum by $Crs^{-3/2}$.
For $E\ge\delta$, Taylor's theorem gives
$R\le(s^2/2)\ee^{-s\delta}(E-\lambda_j)^2$; the exact second moment
pays all remaining initial energies by $Crs^2\ee^{-s\delta}$.
The linear term is smaller. Division by $Z_1\asymp1/(r\sqrt s)$
proves Eq.~\eqref{eq:apcollarexcess}.

There are $O(1/r)$ levels in $\mathcal K$, so its collar
$|\lambda_j-t_0|\le\eta$ has total $d_j=O(r/s)$. Moreover
$t_j\le1+Cr\ee^{s\lambda_j}s^{-3/2}$, whose positive logarithm is
at most $a\eta\mathcal L+o(\mathcal L)+C$. For $t_j\ge1$ use
$q_j(t_j-1)\le d_j$; for the other indices retain the $O(r^2)$
bound. Consequently
\begin{equation}
 \limsup\frac1r\sum_{|\lambda_j-t_0|\le\eta}q_jG(t_j)
                                  \le C_{F,a}\eta.        \label{eq:apentropycollar}
\end{equation}
The constant is independent of small fixed $\eta$. This logarithmic
estimate cannot be replaced by weak convergence alone.

\subsection{Compact bands and activated final energy}

Fix $t_0<c<4$, $M>8$, and $\mathsf P=\one_{[c,M]}(B_\ell)$.
If $a>1$, choose $t_0<e_0<\min(c,1)$; Eq.~\eqref{eq:aptail}
makes the full initial mass above $e_0$ be $o(r/s)$. Below it,
Eq.~\eqref{eq:apbandidentity}, a reduced resolvent analytic at zero,
and Eq.~\eqref{eq:apmom} give
\begin{equation}
 \begin{gathered}
 \left\|\mathsf P\rho\mathsf P-\frac r s\ket{v_\ell}\bra{v_\ell}
                \right\|_1=o(r/s),\qquad\\
 v_\ell=\frac3{m_1}\mathsf P B_\ell^{-1}(Ff_{1,0,\ell}).
 \end{gathered}                                         \label{eq:aprankband}
\end{equation}
The rest of the subthreshold inputs is exponentially paid by
Eq.~\eqref{eq:apcollar}. Strong resolvent convergence and absence of
source atoms imply
\begin{equation}
 \norm{v_\ell}^2\to\nu_F([c,M]),\qquad
 \ip{v_\ell}{B_\ell v_\ell}\to\int_{[c,M]}E\nu_F(\dd E). \label{eq:apbandmoments}
\end{equation}
There are $O_M(1/r)$ band levels. Each final normalized eigenfunction
has left sine coefficient $O(\sqrt r)$ and, after compact transfer,
a right piece bounded by $C\sqrt r(1+x-d)$. This includes energy
four: below it the right piece decays, at it it is linear, and above
it $|\sin(kx)|/k\le x$. Exponential decay of the zero-energy source
then gives $|\ip{\phi_j}{v_\ell}|^2\le C_{F,c,M}r$.
The normalized source energy distribution has entropy
$\mathcal L+O(1)$, between its min-entropy and log-rank bounds.
Measurement and finite-alphabet entropy continuity transfer this
to the actual conditional band distribution with error $o(\mathcal L)$.

If $m=\sum_{[c,M]}p_j$, then $m\sim(r/s)\nu_F([c,M])$ and
$-\log m=\mathcal L+o(\mathcal L)$. Thus
\begin{equation}
 -\sum_{[c,M]}p_j\log p_j=m[2\mathcal L+o(\mathcal L)].   \label{eq:aptwologs}
\end{equation}
One logarithm is the conditional energy entropy and the other the
rare total mass. With $\log Z_4=\mathcal L+o(\mathcal L)$ this gives
\begin{equation}
 \frac1r\sum_{[c,M]}q_jG(t_j)
            \to\int_{[c,M]}(E-a^{-1})\nu_F(\dd E).       \label{eq:apbandentropy}
\end{equation}
The band Gibbs mass is $o(r/s)$, and replacing its KL sum by the
nonnegative $G$ sum changes it by only $O(r/s)=o(r)$.

For $a=1$ choose $1<c<5/4$. The same argument includes all initial
$E<1$; keep
$G_\ell=Z_1^{-1}\mathsf P\one_{[1,\infty)}(A_\ell)
\ee^{-sA_\ell}\mathsf P$ separately. Activated input energy
concentrates at one by Eqs.~\eqref{eq:apheat}--\eqref{eq:aptail},
and its mean squared mass away from the right plateau vanishes by
Eq.~\eqref{eq:aplocal}. The exact identity
$(B_\ell-4)u_n=(E_n-1)u_n-3(1-F)u_n$ therefore concentrates final
energy at four. The fixed band retains all leading activated mass,
and
\begin{align}
 \Tr G_\ell&\sim\ee^{-s},\qquad
 0\le G_\ell\le(\ee^{-s}/Z_1)\mathsf P,\nonumber\\
 (s/r)\Tr G_\ell&\to x,\qquad
 (s/r)\Tr B_\ell G_\ell\to4x.                            \label{eq:apfour}
\end{align}
The positive comparator $(r/s)\ket{v_\ell}\bra{v_\ell}+G_\ell$
approximates the band covariance in trace norm by $o(r/s)$.
Its normalized measured source weights are $O(r)$ and its background
weights are $O(r\sqrt s)$. Its conditional entropy is consequently
still $\mathcal L+o(\mathcal L)$; no source/background orthogonality
is required. The mass coefficient is $\nu_F([c,M])+x$ and the energy
coefficient is $\int_{[c,M]}E\nu_F(\dd E)+4x$. Repeating
Eq.~\eqref{eq:aptwologs} gives
\begin{equation}
 \frac1r\sum_{[c,M]}q_jG(t_j)
              \to\int_{[c,M]}(E-1)\nu_F(\dd E)+3x.       \label{eq:apactivatedentropy}
\end{equation}

\subsection{The entire final-energy tail}

For initial $E<1$ and $M\ge8$, spectral calculus and the exact
overlap identity give
\begin{equation}
 \begin{gathered}
 \ip{u_E}{B\one_{[M,\infty)}(B)u_E}\\
 =9\ip{Fu_E}{\one_{[M,\infty)}(B)B(B-E)^{-2}Fu_E}\\
 \le CM^{-1}\norm{Fu_E}^2.
 \end{gathered}                                         \label{eq:apenergytail}
\end{equation}
For $1\le E\le2$ use Eq.~\eqref{eq:apshiftthree} instead, with
$1-F$ replacing $F$. Inputs above two have prepared mean final
energy $O(\ee^{-2s})$, by heat-weight splitting,
Eq.~\eqref{eq:aptail}, and $B=A+3F$. Thus
\begin{equation}
 \begin{gathered}
 \sum_{\lambda_j\ge M}\lambda_jp_j\\
 \begin{aligned}
 &\le\frac{C_Fr}{sM}
 +\frac{C_F\ee^{-s}}M(s^{-1/2}+r\sqrt s+\ee^{-s})\\
 &\quad+C_F\ee^{-2s}.
 \end{aligned}
 \end{gathered}                                         \label{eq:apfullenergytail}
\end{equation}
For $t_j\ge1$, $q_jG(t_j)\le p_j(s\lambda_j+\log Z_4)$;
the other indices cost $O(r^2)$. Tail probability is at most tail
energy divided by $M$. Bounded $\Xi$ now gives
\begin{equation}
 \limsup\frac1r\sum_{\lambda_j\ge M}q_jG(t_j)
                                      \le C_{F,a}/M.     \label{eq:apentropytail}
\end{equation}
Partition the final spectrum below $t_0-\eta$, in its collar,
between $t_0+\eta$ and $M$, and above $M$. All summands are
nonnegative. Take the device limit first, then $\eta\downarrow0$,
then $M\to\infty$. Equations~\eqref{eq:aplowentropy},
\eqref{eq:apentropycollar}, \eqref{eq:apbandentropy} or
\eqref{eq:apactivatedentropy}, and \eqref{eq:apentropytail} prove
$P/r\to\int(E-a^{-1})_+\nu_F(\dd E)+3x$. The finite first moment
permits exhaustion. Subtracting this result from
Eq.~\eqref{eq:aptotalentropy} proves Eq.~\eqref{eq:thermolaw}.

\section{The sharp-wall infrared work contribution}
\label{app:infrared}

We prove Eqs.~\eqref{eq:abscap}--\eqref{eq:negative} under the
basic regime with $\Xi=O(1)$. In this appendix $F=\one_{x>0}$,
$N=\ell/\sqrt s$, and $O_{nj}=\ip{\phi_j}{u_n}$.

\subsection{Exact low-level overlaps}

The actual subbarrier Sturm labels obey
\begin{align}
 \ell k_n+\delta_1(k_n)&=n\pi,\qquad
 \ell q_j+\delta_4(q_j)=j\pi,\nonumber\\
 \delta_h(k)&=\arctan\frac{k\tanh(\ell\sqrt{h-k^2})}
                                  {\sqrt{h-k^2}}.        \label{eq:apsharpphase}
\end{align}
On fixed small compact energy intervals the normalized left pieces
are $a_h(k)\sin(k(x+\ell))$, with
$a_h(k)=\sqrt{2r}[1+O(r)]$ and
$\sin\delta_h(k)=(k/\sqrt h)[1+O(\ee^{-b_0\ell})]$.
The right pieces are the interface values times
$\sinh(\sqrt{h-k^2}(\ell-x))/\sinh(\ell\sqrt{h-k^2})$.
Their cross integral equals
$[\sqrt{1-k^2}+\sqrt{4-q^2}]^{-1}+O(\ell\ee^{-b_0\ell})$.
Hence, for $j\ne n$,
\begin{align}
 |q_j^2-k_n^2|\,|O_{nj}|^2
 &=\frac{r^2k_n^2q_j^2}{|q_j^2-k_n^2|}\nonumber\\
 &\quad[1+O(k_n^2+q_j^2+r+\ell\ee^{-b_0\ell})].          \label{eq:apoffdiag}
\end{align}
The numerator before dividing by the gap is
$3\ip{\phi_j}{Fu_n}=\pm rk_nq_j[1+O(k_n^2+q_j^2+r+
\ell\ee^{-b_0\ell})]$: the leading factor is
$3(2r)k_n(q_j/2)/(1+2)$.

Put $m=j-n$. Subtracting the exact phases gives
\begin{align}
 (\ell+\delta_4'(\xi))(q_j-k_n)&=\pi m+d(k_n),\nonumber\\
 d(k)&=\delta_1(k)-\delta_4(k)
                \nonumber\\
 &=k/2+O(k^3+k\ee^{-b_0\ell}).                            \label{eq:apgap}
\end{align}
For $m\ne0$ and $k/2\le q_j\le2k$, this implies
$|q_j-k|^{-1}=\ell(\pi|m|)^{-1}[1+O(k/|m|+r)]$.
Since $q^2/(q+k)=k/2+O(|q-k|)$, each near-index term in
Eq.~\eqref{eq:apoffdiag} has leading value $rk^3/(2\pi|m|)$.
The two index ranges have sizes $\ell k/\pi+O(1)$ and
$\ell k/(2\pi)+O(1)$, so their harmonic sums total
$2\log(\ell k)+O(1)$. An empty lower range for the first few modes
is harmless: then $\ell k$ is bounded and is at least $\pi/2$.

For all sufficiently small $0<\delta<\delta_0$ and $k^2<\delta/16$,
the uniform result is
\begin{equation}
 \begin{gathered}
 \sum_{j\ne n,\ q_j^2<\delta}|q_j^2-k^2||O_{nj}|^2\\
   =\frac{rk^3}{\pi}\log(\ell k)+R_{n,\delta},\\
 \begin{aligned}
 |R_{n,\delta}|&\le Cr\{k^2\sqrt\delta+k^3
                  +k^5\log(2+\ell k)\\
 &\quad+(r+\ell\ee^{-b_0\ell})k^3\log(2+\ell k)\}.
 \end{aligned}
 \end{gathered}                                         \label{eq:apharmonic}
\end{equation}
The phase correction costs $O(rk^4\sum m^{-2})$. Replacing
$q^2/(q+k)$ costs $O(r^2k^2(\ell k+1))=O(rk^3)$, because
$k\ge\pi r/2$. The relative $k^2$ correction multiplies the
harmonic sum and produces the $k^5$ term. In the far ranges,
$q<k/2$ costs $O(rk^3)$ since $q^2/|q^2-k^2|\le Cq^2/k^2$;
$2k<q<\sqrt\delta$ costs $O(rk^2\sqrt\delta)$ by counting levels.
This proves Eq.~\eqref{eq:apharmonic} without integrating across
a discrete spectral pole.

\subsection{Every thermal sector and additive error control}

For every fixed $d_0<1$, the complete conditional uncapped mean obeys
\begin{equation}
 \begin{gathered}
 \sum_j|\lambda_j-E_n||O_{nj}|^2
 \le C_{d_0}r[k_n^2+k_n^3\log(2+\ell k_n)]\\
                                      \quad(E_n\le d_0).
 \end{gathered}                                         \label{eq:apworkbound}
\end{equation}
Below a fixed final cutoff less than four, use the numerator bound
$Crkq$, the spacing
$|q_j-k_n|\ge\pi r(|j-n|-1/2)$, the same near/far split, and the
paired shift $O(rE_n)$. Above that cutoff the gap is separated from
zero; $|\omega|\le C\omega^2$ and
\begin{equation}
 \sum_j(\lambda_j-E_n)^2|O_{nj}|^2=9\norm{Fu_n}^2
                                                               \label{eq:apsecondgap}
\end{equation}
pay its entire infinite range by $CrE_n$.

Equation~\eqref{eq:apcollar} bounds
$\sum_{1/2\le E_n<1}\norm{Fu_n}^2$ uniformly, and $N_A(1)=O(\ell)$.
By Cauchy--Schwarz, the conditional absolute first moment is at most
$3\norm{Fu_n}$. Thus the complete collar $[d_0,1)$, $d_0>1/2$,
costs after Gibbs normalization and scaling at most
\begin{equation}
 Cs^{3/2}\sqrt\ell\ee^{-sd_0}
       \le Cs\ee^{-s(d_0-1/2)}=o(1),                      \label{eq:apworkcollar}
\end{equation}
because bounded $\Xi=s\ell\ee^{-s}$ implies $\ell\le C\ee^s/s$.
For $s^{-1/2}\le E_n<d_0$, Eq.~\eqref{eq:apworkbound} and a
Gaussian heat-weight split give a polynomial in $s$ times
$\ee^{-\sqrt s/2}$. The count's factor $\ell$ cancels against $Z$,
and $\log\ell=O(s)$; this intermediate sector is also $o(1)$.

For activated inputs, Eqs.~\eqref{eq:apheat} and
\eqref{eq:apsharpactivated} imply $T/Z\sim\ee^{-s}$ and
$T_-/T=O(r\sqrt s+s^{-1/2})$. Both $h_c(w)=\min(|w|,c)$ and
$h_\infty(w)=|w|$ are one-Lipschitz. Equation~\eqref{eq:apshiftthree}
and Cauchy--Schwarz therefore give
\begin{equation}
 \begin{aligned}
 &\Big|\sum_{E_n\ge1}\ee^{-sE_n}\sum_jh_c(\lambda_j-E_n)|O_{nj}|^2\\
 &\quad-\min(3,c)T\Big|\le3\sqrt{TT_-}.
 \end{aligned}                                         \label{eq:apactivatedwork}
\end{equation}
Its scaled error is $O(\Xi\sqrt{r\sqrt s+s^{-1/2}})=o(1)$, and
its main contribution is $\min(3,c)\Xi$; for uncapped work replace
the coefficient by three. Every activated input and final energy
is included in this estimate.

Only $E_n<s^{-1/2}$ now remains. Gaussian Stieltjes summation gives
\begin{align}
 s\langle E_n\one_{E_n<s^{-1/2}}\rangle_\rho&\to\tfrac12,
 \nonumber\\
 s\langle E_n^2\one_{E_n<s^{-1/2}}\rangle_\rho&\to0,\nonumber\\
 \frac s\pi\langle k_n^3\log(\ell k_n)
                    \one_{E_n<s^{-1/2}}\rangle_\rho
     &=\frac{\log N}{\pi^{3/2}\sqrt s}+o(1).              \label{eq:aplogmoment}
\end{align}
For the logarithmic moment set $v=\sqrt s\,k$. Its continuum numerator
is $(\ell/\pi)s^{-2}\int_0^\infty v^3\ee^{-v^2}
(\log N+\log v)\dd v$. The first integral is $1/2$, the second
finite. The test's total variation is
$O(s^{-3/2}(1+\log N))$, giving an additive scaled quadrature error
$O(r(1+\log N))=o(1)$. The relative partition error
$O(r\sqrt s+\ee^{-s})$ costs only
$O(r\log N+\ee^{-s}\log N/\sqrt s)=o(1)$ against the possibly
diverging main term. The moving Gaussian tail vanishes because
$\log\ell=O(s)$. Unnormalized fixed-cutoff heat errors carry an
$\ell$ prefactor; their normalized contributions are the errors
asserted to vanish here.

The thermal averages of the remainders in
Eq.~\eqref{eq:apharmonic} obey
\begin{align}
 s\langle k^2\sqrt\delta\rangle&=O(\sqrt\delta),\nonumber\\
 s\langle k^3\rangle&=O(s^{-1/2}),\nonumber\\
 s\langle k^5\log(2+\ell k)\rangle
                  &=O(\log\ell/s^{3/2})+o(1)=o(1).       \label{eq:aplogerrors}
\end{align}
The last remainder also vanishes, using $r\sqrt s\to0$,
$\sqrt s=o(\ell)$ and exponential decay in $\ell$. These estimates
are additive even when $\log N/\sqrt s$ diverges.

\subsection{Paired shift, source matching and closure}

At $m=0$, Eq.~\eqref{eq:apgap} gives
\begin{equation}
 \begin{aligned}
 \lambda_n-E_n&=rE_n[1+O(E_n+r+\ee^{-b_0\ell})],\qquad\\
 |O_{nn}|^2&=1+O(E_n+r).
 \end{aligned}                                         \label{eq:appaired}
\end{equation}
For the second bound choose signs so the left sine pieces agree:
their normalized difference has norm
$O(r+\ell|q_n-k_n|)=O(r+k_n)$; their right norms are
$O(\sqrt r\,k_n)$. The paired contribution consequently tends to
$1/2$ by Eq.~\eqref{eq:aplogmoment}.

Fix $0<\delta<\min(\delta_0,c)$. The remaining final modes below
$\delta$ contribute $\log N/(\pi^{3/2}\sqrt s)+R_\delta$, with
$\limsup|R_\delta|\le C\sqrt\delta$. For all final energies above
$\delta$ use the bounded reduced spectral function
\begin{equation}
 G_{E,\delta}(z)=\one_{z\ge\delta}\frac{h_c(z-E)}{(z-E)^2}.
                                                               \label{eq:apreducedtest}
\end{equation}
For $E<\delta/2$ it is uniformly bounded and differs from
$G_{0,\delta}$ by $O_\delta(E)$ in supremum norm. For uncapped work
it is simply $\one_{z\ge\delta}/(z-E)$. Exact matching gives
$Fu_n/(\sqrt{2r}k_n)\to f$ in $L^2$, uniformly over the moving
low-input window. The overlap identity implies
\begin{align}
 \sum_{\lambda_j\ge\delta}h_c(\lambda_j-E_n)|O_{nj}|^2
 &=9\ip{Fu_n}{G_{E_n,\delta}(B_\ell)Fu_n}\nonumber\\
 &=2rE_n[\Phi_\delta(c)+o(1)],\nonumber\\
 \Phi_\delta(c)&=\int_{[\delta,\infty)}\min(E,c)\nu(\dd E).
                                                               \label{eq:apsourcematch}
\end{align}
Strong resolvent convergence applies at the cutoff because the
source has no atom there; the bounded function includes the whole
infinite final range. Thermal averaging gives $\Phi_\delta(c)$.
The omitted source first moment is $O(\sqrt\delta)$. Combine these
three disjoint final sectors and the preceding input bounds, take
the device limit, then send $\delta\downarrow0$. This proves
Eq.~\eqref{eq:abscap}. The same proof, directly with the uncapped
test and $\int E\nu(\dd E)=1$, proves Eq.~\eqref{eq:absolute}.

The exact signed mean is $3\Tr\rho F$. Its low and activated
contributions are $3r/(2s)$ and $3\ee^{-s}$, respectively, proving
the signed formula in Eq.~\eqref{eq:negative}. Subtract it from
the absolute-work result. As an independent sign check, the negative
near-index sum uses just $m<0$, half the harmonic coefficient; far
negative terms cost $O(rk^3)$, and the activated negative mean is
$o(1)$ after scaling since $(-w)_+$ is one-Lipschitz and vanishes
at three. Finally, when $s/\mathcal L\to a>1$,
$\log N/\sqrt s\sim\sqrt s/a$ and $\Xi\to0$, which proves both
ratios in Eq.~\eqref{eq:ratio}.

\section{Source reconstruction and frequency rejection}
\label{app:reconstruction}

\subsection{Positive transform, finite window and frequency rejection}

For a positive measure on $(0,\infty)$ with first moment $M<\infty$,
$C(t)$ and $\Phi(c)$ are finite and $|C'(t)|\le M/2$ by dominated
differentiation. Integration by parts and the Dirichlet sine integral
give $\int_0^\infty(1-\cos ut)t^{-2}\dd t=\pi|u|/2$.
Product-to-sum consequently gives
\begin{equation}
 \int_0^\infty\frac{(1-\cos ct)(1-\cos Et)}{t^2}\dd t
                          =\frac\pi2\min(c,E).            \label{eq:apcosine}
\end{equation}
Tonelli applied to the nonnegative integrand proves
Eq.~\eqref{eq:positive}. Its kernel has mass $2c$, so the
uniform-additive-error constant $2c$ is exact.

If $\nu'(E)\le K E^{-3/2}$ for $0<E<E_*$, then
$C(t)\le K\sqrt{\pi/2}\sqrt t+M/E_*$. Therefore the positive
omitted time tail satisfies
\begin{equation}
 0\le\Phi(c)-\Phi_T(c)
 \le\frac{16K}{\sqrt{2\pi T}}+\frac{8M}{\pi E_*T}.      \label{eq:aptimetail}
\end{equation}
For exact midpoint weights on a partition of maximum mesh $\Delta$,
source sample errors at most $\varepsilon$ give
\begin{equation}
 \left|\sum_iw_i\widehat C(t_i)-\Phi(c)\right|
 \le\frac{16K}{\sqrt{2\pi T}}+\frac{8M}{\pi E_*T}
                  +2c\varepsilon+\frac{cM\Delta}{2}.     \label{eq:apquadrature}
\end{equation}
The last term uses $|C'|\le M/2$, a maximal midpoint distance
$\Delta/2$, and $\sum w_i\le2c$.

These constants can be explicit for all fixed-width profiles.
Let $n_d=2d+1/2$, $m_*=(2d+1)^{-1}$ and
\begin{equation}
 E_*=\min\{1/4,\gamma_d/4,m_*/(8n_d)\},\qquad
 K=\frac{144n_d^2}{\pi m_*^4}.                            \label{eq:approfileconstants}
\end{equation}
Equation~\eqref{eq:apode} implies
$\norm{f_{4,E}}\le2\sqrt{n_d}$ and $m_4(E)\ge m_*/2$
on this interval; also $\norm g\le\sqrt{n_d}$. Substitution into
Eq.~\eqref{eq:apdensity} proves the stated threshold bound.
The zero-energy energy identity gives $\dot m_1\le m_1$, hence
$M_F\le3/m_*$.

For a sharper subsequent tail, if
$|\nu'(E)-kE^{-3/2}|\le BE^{-1/2}$ below $E_*$, the signed
difference from $kE^{-3/2}\dd E$ on the entire positive axis has
total variation at most
$R_0=2B\sqrt{E_*}+M/E_*+2k/\sqrt{E_*}$. Thus
$C(t)=A\sqrt t+R(t)$, $A=k\sqrt{\pi/2}$, $|R(t)|\le R_0$.
An integration by parts of the cosine tail yields
\begin{equation}
 \left|\Phi(c)-\Phi_T(c)-\frac{8A}{\pi\sqrt T}\right|
 \le\frac{8A}{\pi cT^{3/2}}+\frac{8R_0}{\pi T}.           \label{eq:apleadingtail}
\end{equation}
For the actual source $A=D_F^2/\sqrt{2\pi}$; the remainder is
not uniform as the cap tends to zero.

Exact all-time source data determine all caps, then the measure via
$\Phi'_+(c)=\nu((c,\infty))$ and
$\Phi'_-(c)=\nu([c,\infty))$, by dominated one-sided difference
quotients away from zero. This is not a stable inversion in a strong
measure norm. In particular adding $\varepsilon\delta_H$ changes
$C$ uniformly by at most $\varepsilon$ but changes its first moment
by $\varepsilon H$. Resonance aliasing in the main text therefore
does not contradict all-time injectivity. It can even preserve two
moments with smooth densities: take a smooth probability density
$\psi$ supported in $(0,c)$, of mean $c/2$, and a common base
$kE^{-3/2}\ee^{-E}\dd E$. Adding respectively
$\beta(3T_3\psi+T_9\psi)/4$ and
$\beta(\psi+3T_6\psi)/4$, where $T_j$ translates by $j$, gives
equal first and second moments and identical resonant samples, but
caps differing by $\beta c/8$. These are measure-class examples,
not an asserted pair of realizable wall profiles.

For the actual finite work probability, Eq.~\eqref{eq:apcosine}
also proves the exact all-time identity stated after
Eq.~\eqref{eq:positive}. For fixed $0<\eta<c$ put
$h_{\eta,c}(w)=\min(|w|,c)-\min(|w|,\eta)$.
Its quotient by $w^2$, set to zero at zero, is bounded continuous
with norm at most $1/(4\eta)$. Equation~\eqref{eq:weightedlaw}
therefore proves Eq.~\eqref{eq:reject}. Splitting the source cap at
$\eta$ and $E_*$ gives the cutoff budget
\begin{equation}
 0\le\Phi_F(\eta)\le4K\sqrt\eta+\eta M_F/E_*
                               \quad(\eta\le E_*).        \label{eq:apcutoffbudget}
\end{equation}
Subtracting two absolutely integrable cosine identities gives its
signed time kernel. Its total absolute weight is at most
$2(c+\eta)$; since $p_{r,s}\le1$, the scaled actual-device tail
beyond $T$ is at most $8(s/r)/(\pi T)$. The ordered limiting
filter does not supply an uncharged finite observation protocol.

% Standalone scientific appendices; insert before the bibliography.
% Uses only macros already defined in the root manuscript preamble.
\section{Resonant source selection and complete entropy bounds}
\label{app:control}

We prove the actual conditional-state assertion, including its
infinite-dimensional entropy limit. All operators in this appendix
are in the boundary coordinate on $(-\ell,\ell)$, with
$\ell=1/r$. The resonance $\tau=2\pi n/3$ is fixed, where $n$ is a
positive integer. Constants may depend on this resonance and on a
fixed smooth profile, but not on $\ell$ or $s$.
For a density operator $\sigma$ and $q>0$, $q\ne1$, the
R\'enyi entropy is $S_q(\sigma)=\log\Tr\sigma^q/(1-q)$;
the order-one entropy is $S(\sigma)=-\Tr\sigma\log\sigma$.

\subsection{The sharp low-energy covariance}

For the sharp wall, put $Q=\one_{(0,\ell)}$ and
$A=-\partial_x^2+Q$, $B=-\partial_x^2+4Q$. At $E=k^2<1$ define
\begin{align}
 \kappa&=\sqrt{1-E},\qquad d(E)=\kappa\coth(\kappa\ell),\nonumber\\
 f_{E,\ell}(x)&=\frac{\sinh[\kappa(\ell-x)]}
                         {\sinh(\kappa\ell)}\one_{(0,\ell)}(x).
                                                        \label{eq:ctrlf}
\end{align}
Matching the left sine and right hyperbolic sine gives
$k\ell+\arctan[k/d(k^2)]=j\pi$. If $u_E$ is normalized and
$b_E=u_E(0)$, direct integration gives
\begin{align}
 Qu_E&=b_E f_{E,\ell},\qquad |b_E|^2=\frac{2E}{D_\ell(E)},\nonumber\\
 D_\ell(E)&=\ell+[\ell-2\ell E]\operatorname{csch}^2(\kappa\ell)
                       +\frac{\coth(\kappa\ell)}\kappa.
                                                        \label{eq:ctrlamp}
\end{align}
On $0\le E\le1/2$, one has
$\ell\le D_\ell(E)\le\ell+3$ for large $\ell$.
The phase is increasing and its count differs from $\ell k/\pi$
by at most one. Consequently, for each fixed integer $j\ge0$,
\begin{equation}
 M_j:=\sum_{E<1/2}|b_E|^2\ee^{-sE}E^j
 =\frac{\Gamma(j+3/2)}\pi s^{-j-3/2}[1+o(1)].
                                                        \label{eq:ctrlmom}
\end{equation}
Indeed, Stieltjes integration of $k^{2j+2}\ee^{-sk^2}$ against
the count gives discrepancy at most twice its maximum. After
multiplication by $2/\ell$, its relative size is
$O_j(\sqrt s/\ell)$. The normalization replacement costs
$O(\ell^{-1})$, and the omitted Gaussian integral above
$k=1/\sqrt2$ is exponentially small. These errors vanish in
Eq.~\eqref{eq:regime}. Write
$m_s=(2\sqrt\pi s^{3/2})^{-1}$, so $M_0\sim m_s$.

Duhamel's formula gives the exact column identity
\begin{equation}
 K_tu_E=b_E\ee^{-itE}V_\ell(E),\qquad
 V_\ell(E)=\frac32h_t(B-E)f_{E,\ell}.
                                                        \label{eq:ctrlcolumn}
\end{equation}
The low accepted covariance is thus
\begin{equation}
 S=\sum_{E<1/2}|b_E|^2\ee^{-sE}
                 \ket{V_\ell(E)}\bra{V_\ell(E)}.
                                                        \label{eq:ctrlcov}
\end{equation}
Different initial eigenstates contribute positive outer products;
their amplitudes are not coherently averaged.

Let $T_\ell$ be the Dirichlet inverse of $-\partial_x^2+1$ on
$(0,\ell)$. The inhomogeneous right boundary problem gives
$f_{E,\ell}=(I-ET_\ell)^{-1}f_{0,\ell}$, with
$\norm{T_\ell}\le1$ and $\norm{f_{0,\ell}}\le1/\sqrt2$.
Using $h_t(B-E)=-i\int_0^t\ee^{-iuB}\ee^{iuE}\dd u$, the
Taylor coefficients of $V_\ell(E)=\sum_{j\ge0}E^jV_{j,\ell}$ obey
\begin{equation}
 \norm{V_{j,\ell}}\le M_t:=\frac3{2\sqrt2}\int_0^t\ee^u\dd u.
                                                        \label{eq:ctrltaylor}
\end{equation}
Therefore
$\norm{S-M_0\ket{V_\ell(0)}\bra{V_\ell(0)}}_1\le C_tM_1$.
Increasing-interval Dirichlet form convergence implies strong
resolvent convergence to $B_\infty=-\partial_x^2+4\one_{x>0}$
and hence strong propagator convergence at every fixed time.
Since $f_{0,\ell}\to f=\ee^{-x}\one_{x>0}$ in $L^2$,
\begin{equation}
 V_\ell(0)\longrightarrow v_t=\frac32h_t(B_\infty)f,
 \qquad\norm{v_t}^2=C(t)>0.                              \label{eq:ctrlsource}
\end{equation}
For positivity, $B_\infty$ has no eigenvalues: positivity excludes
negative ones, and at nonnegative energy its free left solution is
affine or oscillatory and cannot be square integrable unless zero.
ODE uniqueness then makes the whole solution zero. Its spectral
measure for nonzero $f$ is atom-free, whereas the real zero set of
$h_t$ is countable. This proves the strict inequality in
Eq.~\eqref{eq:ctrlsource}. The actual comparison projection is
$\Pi_{r,t}=\ket{V_\ell(0)}\bra{V_\ell(0)}/\norm{V_\ell(0)}^2$.

\subsection{An all-energy boundary estimate at resonance}

Let $H=-\partial_x^2+V$, $0\le V\le v$, on an interval with a
unit subinterval adjacent to the point under consideration.
For $R=(H+1)^{-1}$, spectral calculus and the differential equation
give
\begin{equation}
 \begin{gathered}
 \norm{Rg}\le\norm g,\quad\\
 \norm{(Rg)'}\le\tfrac12\norm g,\quad\\
 \norm{(Rg)''}\le(v+2)\norm g.                            \label{eq:ctrlres}
 \end{gathered}
\end{equation}
The elementary endpoint inequality
$|h(0)|^2\le2\norm h^2+\norm{h'}^2$ on a unit interval,
applied to $Rg$ and $(Rg)'$, gives by Parseval
\begin{equation}
 \sum_E\frac{|u_E(0)|^2+|u'_E(0)|^2}{(1+E)^2}
 \le D_v:=\frac{11}4+(v+2)^2.                            \label{eq:ctrltrace}
\end{equation}
This also defines the bounded resolvent-smoothed distributions
$R\delta_0,R\delta'_0$. The constant is independent of total
length, local level spacing, and threshold proximity. Separated
remote Neumann conditions give the same estimate.

At resonance the function
\begin{equation}
 q_\tau(z)=\frac{\ee^{-i\tau z}-1}{z(z-3)}                \label{eq:ctrlq}
\end{equation}
is entire. Green's identity on the controlled interval gives, for
every actual initial and final eigenstate,
\begin{align}
 (\lambda-E)\ip{\phi_\lambda}{u_E}
 &=3\ip{\phi_\lambda}{Qu_E},\nonumber\\
 (\lambda-E-3)\ip{\phi_\lambda}{Qu_E}
 &=-\overline{\phi_\lambda(0)}u'_E(0)
              +\overline{\phi'_\lambda(0)}u_E(0).
                                                        \label{eq:ctrlgreen}
\end{align}
Physical outer endpoint terms vanish by the Dirichlet conditions.
Combining the identities, including their removable equal-energy
values, proves
\begin{equation}
 K_\tau u_E=-\frac32\ee^{-i\tau E}q_\tau(B-E)
                    [u'_E(0)\delta_0+u_E(0)\delta'_0].
                                                        \label{eq:ctrlboundary}
\end{equation}
Its distributional meaning is fixed by
$q_\tau(B-E)\delta_0^{(i)}=[q_\tau(B-E)(B+1)]
(B+1)^{-1}\delta_0^{(i)}$, whose bracket is bounded.

Near zero use $|\ee^{-i\tau x}-1|\le\tau|x|$; near three use
the resonant version of that estimate; away from both points use
the numerator bound two. These give
$\sup_{x\in\mathbb R}(1+|x|)|q_\tau(x)|\le5\tau/2$.
Since $1+\lambda\le(1+E)(1+|\lambda-E|)$, the input and output
versions of Eq.~\eqref{eq:ctrltrace} imply, for each $d_0>0$,
\begin{align}
 &\Tr K_\tau\one_{[d_0,\infty)}(A)\ee^{-sA}K_\tau^*
 \nonumber\\[-2pt]
 &\hspace{15pt}\le C_\tau\sup_{E\ge d_0}\ee^{-sE}(1+E)^4
 \le C_\tau(1+d_0)^4\ee^{-sd_0},                         \label{eq:ctrlhightrace}
\end{align}
where the last inequality holds for $s\ge4/(1+d_0)$.
Taking $d_0=1/2$ pays the complete complement of
Eq.~\eqref{eq:ctrlcov}, including every threshold and high mode.
No box-volume factor remains.

\subsection{All-index tails and entropy}

The first $m$ Taylor vectors in Eq.~\eqref{eq:ctrltaylor} leave
error at most $2M_\tau E^m$ on $0\le E\le1/2$. Positivity and
the variational characterization of eigenvalue sums yield
\begin{equation}
 \sum_{j>m}\lambda_j(S)\le4M_\tau^2M_{2m}.               \label{eq:ctrllowtail}
\end{equation}
The actual cutoff gives $E^{2m}\le4^{-(m-1)}E^2$, so
Eqs.~\eqref{eq:ctrlmom} and \eqref{eq:ctrlsource} imply
\begin{equation}
 \sum_{j>m}\lambda_j(S/\Tr S)
 \le A_\tau s^{-2}4^{-(m-1)},\qquad m\ge1.               \label{eq:ctrllowgeo}
\end{equation}
Summing $-\lambda\log\lambda$ under this envelope, including the
$O(s^{-2})$ deficit of the leading eigenvalue, proves
$S(S/\Tr S)=O_\tau(s^{-2}\log s)$. For each fixed $0<q<1$ it
also proves $\Tr(S/\Tr S)^q=1+O_{\tau,q}(s^{-2q})$.

Write $G=K_\tau\one_{[1/2,\infty)}(A)\ee^{-sA}K_\tau^*$.
The two vector curves
$g_i(E)=q_\tau(B-E)\delta_0^{(i)}$, $i=0,1$, are norm-analytic
on the strip $|\operatorname{Im}E|\le1$ and satisfy
\begin{equation}
 \norm{g_i(E)}\le A_\tau(1+|E|).                         \label{eq:ctrlstrip}
\end{equation}
Indeed the entire divided difference bounds $q_\tau$ near its
removable points, while its quotient gives quadratic decay at large
real part; hence $(1+|z|)|q_\tau(z)|$ is bounded on the strip.
Factoring through the local resolvent vectors proves
Eq.~\eqref{eq:ctrlstrip}.

In the bin $[1/2+j,3/2+j)$ expand both curves at its midpoint on
the radius-one circle. Retaining $m$ terms costs at most $2m$
dimensions and has vector error $A_\tau(j+3)2^{-m}$.
Equation~\eqref{eq:ctrltrace} bounds the sum of squared input
boundary coefficients in the bin by $D_1(j+3)^2$. Retaining their
Cauchy--Schwarz cross bound, the trace outside the approximation
space is at most
\begin{equation}
 A_\tau(j+3)^4\ee^{-s(1/2+j)}4^{-m}.                    \label{eq:ctrlbin}
\end{equation}
Choose $m_j=M-j$ for $0\le j<M$ and none above. The total
dimension is at most $M(M+1)$. Both the retained-bin errors and
the omitted bins sum because $4\ee^{-s}<1$ for $s\ge2$.
Thus
\begin{align}
 \sum_{k>M(M+1)}\lambda_k(G)&\le A_\tau\ee^{-s/2}4^{-M},\nonumber\\
 \lambda_k(G)&\le A'_\tau\ee^{-s/2}\ee^{-c\sqrt k}.
                                                        \label{eq:ctrlhightail}
\end{align}
This controls both complete input and final spectral tails.

Let $T=\Tr(S+G)\sim C(\tau)m_s$ and $Y=G/T$.
Its eigenvalues are bounded by
$\varepsilon_s\ee^{-c\sqrt k}$ with
$\varepsilon_s=O_\tau(s^{3/2}\ee^{-s/2})$.
The monotonicity of $-x\log x$ below $1/\ee$ and summability
of this envelope show
\begin{equation}
 \begin{gathered}
 S(Y)=O_\tau(s^{5/2}\ee^{-s/2}),\qquad\\
 \Tr Y^q=O_{\tau,q}(\varepsilon_s^q),\quad 0<q<1.
                                                        \label{eq:ctrlhighentropy}
 \end{gathered}
\end{equation}
Here $S(X)=-\Tr X\log X$ is also used for subnormalized positive
operators. Positive-sum entropy bounds give
\begin{equation}
 S(\sigma_\tau)=O_\tau(s^{-2}\log s+s^{5/2}\ee^{-s/2})\to0.
                                                        \label{eq:ctrlentropy}
\end{equation}
For $0<q<1$ use trace-power subadditivity; for $q>1$ the largest
eigenvalue tending to one suffices. This proves every fixed positive
R\'enyi limit. Finally,
$Z=(1+\ee^{-s})\ell/(2\sqrt\pi\sqrt s)+O(1)$ by
Dirichlet/Neumann splitting. Division of the accepted trace by
this complete partition proves $p_\tau\sim C(\tau)r/s$.

\subsection{Fixed smooth plateau profiles}

We supply the complete smooth high-sector argument; smooth entropy
is not inferred from the sharp trace bound. Suppose $F=0$ to the
left of $-d$ and $F=1$ to the right of $d$, as in the main text.
The right Dirichlet gap, boundary normalization, and moments are
proved in Eqs.~\eqref{eq:apode}--\eqref{eq:apmom}.
For explicit analytic radii one can take
\begin{equation}
 \delta_F=\min\left\{\frac14,\frac{\gamma_d}4,
                   \frac1{8(2d+1)(2d+1/2)}\right\}>0.
                                                        \label{eq:ctrldelta}
\end{equation}
The zero-energy right solution satisfies
$\norm{f_{1,0,\ell}}^2\le2d+1/2$, and the resolvent series
gives $\norm{f_{1,E,\ell}}\le2\sqrt{2d+1/2}$ for
$|E|\le2\delta_F$.
The exact source curve is
\begin{equation}
 V_\ell(E)=\frac32h_\tau(B-E)Ff_{1,E,\ell}.              \label{eq:ctrlsmoothsource}
\end{equation}
It is uniformly norm-analytic on that disk. Cauchy's estimate
on radius $2\delta_F$, evaluated for $E\le\delta_F$, gives the
same normalized geometric low tail as Eq.~\eqref{eq:ctrllowgeo}.
The boundary moment has the factor $m_1^{-2}$, so its accepted
trace is asymptotic to $C_F(\tau)m_s$.

For the complete high sector split $F=F_c+Q_d$, where
$F_c=F\one_{(-d,d)}$ and $Q_d=\one_{(d,\ell)}$.
An exact all-energy identity is
\begin{align}
 K_\tau u_E=\frac32\ee^{-i\tau E}\bigl\{
 &h_\tau(B-E)F_cu_E\nonumber\\[-2pt]
 &-q_\tau(B-E)[u'_E(d)\delta_d+u_E(d)\delta'_d]\bigr\}.
                                                        \label{eq:ctrlsmoothhigh}
\end{align}
The interface distributions result from the artificial split, not a
discontinuity of $F$. The plateau identity
$(B-E-3)Q_du_E=-u'_E(d)\delta_d-u_E(d)\delta'_d$, multiplied
by the removable resonant quotient, proves the formula.

On $[-d,d]$ write
$u_E=u_E(d)\phi_E+u'_E(d)\psi_E$, where the two solutions have
terminal data $(1,0)$ and $(0,1)$ at $d$. Their ODE integral
series is entire in $E$, and first-order Gronwall gives a bound
$\exp[2d(2+|E|)]$ on each solution and derivative.
Thus Eq.~\eqref{eq:ctrlsmoothhigh} is $\ee^{-i\tau E}$ times
two norm-analytic vectors with the actual boundary coefficients.
On the radius-one circle centered at $\delta_F+j+1/2$, their
norms are at most $C_{d,\tau}(j+3)\ee^{2dj}$.
The local trace inequality at $d$ is unchanged. A bin with $m$
Taylor terms consequently has trace error at most
\begin{equation}
 C_{d,\tau}(j+3)^4\ee^{-s\delta_F}
                     \ee^{-(s-4d)j}4^{-m}.              \label{eq:ctrlsmoothbin}
\end{equation}
The same allocation $m_j=M-j$ proves, for $s\ge4d+2$,
\begin{equation}
 \sum_{k>M(M+1)}\lambda_k(G_F)
       \le C_{d,\tau}\ee^{-s\delta_F}4^{-M},             \label{eq:ctrlsmoothtail}
\end{equation}
where $G_F$ includes every input energy at least $\delta_F$.
In particular it has the all-index stretched-exponential envelope
needed after conditioning. Repeating the normalization and entropy
argument gives
\begin{equation}
 S(\sigma_\tau)=O_{F,\tau}
       (s^{-2}\log s+s^{5/2}\ee^{-s\delta_F})\to0,        \label{eq:ctrlsmoothentropy}
\end{equation}
and all fixed positive R\'enyi limits.
The positive source limit follows from strong resolvent convergence
and the absence of eigenvalues of the full-line operator, whose
left half-line is free. Its coefficient is precisely
$\norm{(3/2)h_\tau(B_\infty)Ff_1}^2/m_1^2=C_F(\tau)$.
These constants concern a fixed profile in penetration-length
units; the theorem does not keep a positive physical transition
width fixed as $r\to0$.

\section{One controlled wall versus two in the same Gibbs well}
\label{app:switch}

Use the connected initial well in Eq.~\eqref{eq:twowell}.
Restriction to the right half, multiplied by $\sqrt2$, is unitary
from each initial parity subspace. In right-interface coordinates
the half interval is $(-a,\ell)$, $a=1/(2r)$,
$\ell=1/r$. Even parity gives a Neumann condition at $-a$;
odd parity gives a Dirichlet condition. This decomposition is an
exact representation of the initial Sturm problem and introduces
no physical midpoint wall.

For $E=k^2<1$, the quantization conditions are
\begin{equation}
 ka+\arctan[k/d(E)]=
 \begin{cases}(j-1/2)\pi,&\text{even},\\j\pi,&\text{odd},\end{cases}
                                                        \label{eq:switchphase}
\end{equation}
with the right function in Eq.~\eqref{eq:ctrlf}. Integration
gives the normalized half-state boundary amplitude
\begin{equation}
 b_E^2=\frac{2E}
 {a+[a-(a+\ell)E]\operatorname{csch}^2(\kappa\ell)
                      +\coth(\kappa\ell)/\kappa}.       \label{eq:switchamp}
\end{equation}
The denominator lies between $a$ and $a+3$ for $E\le1/2$ and
large $a$. Each parity count differs from $ak/\pi$ by at most
one. The proof of Eq.~\eqref{eq:ctrlmom} therefore gives
\begin{equation}
 M_{j,+}\sim M_{j,-}\sim
       \frac{\Gamma(j+3/2)}\pi s^{-j-3/2}.               \label{eq:switchmom}
\end{equation}
Equal leading weights follow from this computation, not reflection
symmetry alone.

When only the right wall is controlled, the actual full-state
boundary amplitude is $\beta_E=b_E/\sqrt2$. Its union moment is
$(M_{j,+}+M_{j,-})/2$, with the single-wall leading coefficient.
The final operator is on $(-2\ell,\ell)$ with height one on
$(-2\ell,-\ell)$ and height four on $(0,\ell)$. It does not
preserve initial parity. For both initial parity sectors,
\begin{equation}
 K_Ru_E=\beta_E\ee^{-i\tau E}\frac32
                 h_\tau(B_R-E)f_{E,\ell}.               \label{eq:switchcolumn}
\end{equation}
Thus there is one common source curve. Both initial parity
populations approach the same output vector. The uncontrolled
left barrier and the endpoints recede in the source limit, giving
the same full-line $v_\tau$ as Eq.~\eqref{eq:ctrlsource}.
Only the controlled interface contributes to
Eq.~\eqref{eq:ctrlboundary}; the uncontrolled left wall commutes
with $Q_R$. The preceding local trace, analytic-curve, and entropy
bounds apply unchanged. Hence the accepted trace is
$C(\tau)m_s[1+o(1)]$ and the output entropy tends to zero.

When both walls are controlled, both final parity blocks are
exactly invariant and orthogonal. Their half-interval operators
are the same sharp $A,B$ pair; only the remote midpoint condition
differs. At a negative spectral parameter the reflected Green
terms from a remote Neumann or Dirichlet endpoint decay
exponentially, giving the same strong source limit for both.
Local resolvent estimates and their integrations by parts hold
with either remote condition. Each block therefore has accepted
trace $C(\tau)m_s[1+o(1)]$ and normalized entropy tending to zero.
Their weights approach $1/2$. If the weight is $z$, the exact
identity
\begin{equation}
 S(z\sigma_+\oplus(1-z)\sigma_-)
 =h_2(z)+zS(\sigma_+)+(1-z)S(\sigma_-)                   \label{eq:switchentropy}
\end{equation}
where $h_2(z)=-z\log z-(1-z)\log(1-z)$, proves the limit
$\log2$. The block trace-power identity proves
the same limit for each fixed positive R\'enyi order.

For the unchanged input, three-interval bracketing gives
\begin{equation}
 Z=\frac{1+2\ee^{-s}}{2\sqrt\pi r\sqrt s}+O(1).        \label{eq:switchpartition}
\end{equation}
It proves the two stated probabilities and their ratio two,
as well as input purity $\sim\sqrt{2\pi}r\sqrt s$.
There is no finite-parameter identity $K_2=K_L+K_R$, and the
calculation does not assert that arbitrary interface counts
equal accepted rank.

\section{Survival of repeated noncommuting heralds}
\label{app:repeat}

Here the resonance is the fundamental one, $\tau=2\pi/3$.
We prove injectivity for the actual full-line difference of
propagators and then transfer the complete thermal estimates.
This is needed because the difference of two unitaries can be
nonnormal and can even have a zero square; a scalar spectral
filter argument does not apply to $K^m$.

\subsection{The sharp step}

Suppose $Kg=0$, $\norm g=1$. Set $c=3/2$,
$W(u)=\ee^{iuA}\ee^{-iuB}$,
$Z(u)=\ee^{icu}W(u)$, and
$S(u)=\ee^{iuA}(2Q-I)\ee^{-iuA}$.
Bounded perturbation makes $Z(u)g$ strongly continuously
differentiable, with
\begin{equation}
 Z'(u)=-icS(u)Z(u),\qquad S(u)^*=S(u),\quad S(u)^2=I.
                                                        \label{eq:repeatcurve}
\end{equation}
On the real unit sphere of the complex Hilbert space, the curve
$\psi(u)=Z(u)g$ has speed $c$, joins $g$ to $-g$, and has length
$c\tau=\pi$. It is a shortest great semicircle. Explicitly, the
two spherical distances from $\psi(u)$ to the antipodes sum to
$\pi$ and are at most $cu$ and $\pi-cu$; both bounds are equalities.
Writing $\psi(u)=\cos(cu)g+\sin(cu)n(u)$ with $n(u)$ a real
orthogonal unit vector gives
$\norm{\psi'}^2=c^2+\sin^2(cu)\norm{n'}^2$, so $n$ is
constant and fixed by the initial derivative. Consequently
\begin{equation}
 W(u)g=(I-Q)g+\ee^{-3iu}Qg.                              \label{eq:repeatgeodesic}
\end{equation}
The real-sphere argument retains the Ramsey relative phase;
projective Hilbert space would identify the two endpoints and
would not give this conclusion.

Inserting Eq.~\eqref{eq:repeatgeodesic} into the interaction
equation gives
\begin{equation}
 Q\ee^{-iuA}(I-Q)g
       =\ee^{-3iu}(I-Q)\ee^{-iuA}Qg.                    \label{eq:repeatsupport}
\end{equation}
The two sides have disjoint supports and therefore both vanish
for $0<u<\tau$. Each initial half-line piece would stay confined
to its original half-line.

Persistent confinement is excluded by semibounded time analyticity,
the principle underlying the positive-observable alternative in
Ref.~\cite{Hegerfeldt1998}; we give its application directly here.
For $f_-=(I-Q)g$ and any test vector $k$ supported in $x>0$,
$F_k(z)=\ip{k}{\ee^{-izA}f_-}$ is bounded analytic for
$\operatorname{Im}z<0$, continuous on the real boundary, and
zero on an open interval. Schwarz reflection across a smaller
interval with zero boundary values followed by the identity
theorem makes it identically zero. The heat evolution and all
resolvents $(A+\lambda)^{-1}f_-$ therefore vanish on $x>0$.
Direct matching gives the cross-interface resolvent kernel
\begin{equation}
 \begin{gathered}
 G_\lambda(x,y)=\frac{\ee^{-qx}\ee^{py}}{p+q},\quad\\
 x>0>y,\quad p=\sqrt\lambda,\quad q=\sqrt{\lambda+1}.
                                                        \label{eq:repeatresolvent}
 \end{gathered}
\end{equation}
Thus $\int_{-\infty}^0\ee^{py}f_-(y)\dd y=0$ for every
$p>0$. Analytic continuation to $p=a+i\xi$ and Fourier
uniqueness for the exponentially damped half-line function
prove $f_-=0$. The reversed cross kernel gives the same conclusion
for $f_+=Qg$, initially for $q>1$ and then by analytic continuation.
This contradiction proves $\ker K=\{0\}$.

\subsection{Smooth plateau and strictly interior profiles}

For a smooth nondecreasing plateau profile replace $2Q-I$ by
$2F-I$ in Eq.~\eqref{eq:repeatcurve}. Its norm is at most one,
so the curve from $g$ to $-g$ has length at most $\pi$. Equality
is necessary, and every speed equals $c$. It follows that
\begin{equation}
 4\ip{\ee^{-iuB}g}{F(1-F)\ee^{-iuB}g}=0,
                  \qquad0\le u\le\tau.                 \label{eq:repeatsmoothzero}
\end{equation}
Let $a=\sup\{x:F(x)=0\}$ and $b=\inf\{x:F(x)=1\}$.
Continuity and monotonicity give $a<b$ and $0<F<1$ on $(a,b)$.
Equation~\eqref{eq:repeatsmoothzero} makes the evolution vanish
there, including at zero time. The preceding time analyticity
argument implies that $(B+\lambda)^{-1}g$ vanishes there for
every $\lambda>0$.

Let $\phi_-$ and $\phi_+$ solve $(B+\lambda)\phi=0$ and agree
with $\ee^{\sqrt\lambda(x-a)}$ for $x\le a$ and
$\ee^{-\sqrt{\lambda+4}(x-b)}$ for $x\ge b$, respectively.
They are linearly independent: proportionality would give a
negative-energy $L^2$ eigenfunction, contradicting $B\ge0$.
Their Wronskian is nonzero, and the Green formula, verified by
its unit derivative jump, gives on $(a,b)$
\begin{equation}
 [(B+\lambda)^{-1}g](x)=c_-(\lambda)\phi_+(x)
                                      +c_+(\lambda)\phi_-(x).
                                                        \label{eq:repeatsmoothgreen}
\end{equation}
Up to the same nonzero Wronskian factor the coefficients are
\begin{align}
 c_-(\lambda)&=\int_{-\infty}^a
                \ee^{\sqrt\lambda(y-a)}g(y)\dd y,\nonumber\\
 c_+(\lambda)&=\int_b^\infty
                \ee^{-\sqrt{\lambda+4}(y-b)}g(y)\dd y.
                                                        \label{eq:repeatlaplace}
\end{align}
There is no source integral inside $(a,b)$ because $g$ vanishes
there. Linear independence forces both coefficients to vanish;
the two Laplace uniqueness arguments prove $g=0$.
This proves injectivity for the fixed smooth plateau family.

For completeness, a multiplier satisfying $0<F<1$ almost
everywhere, such as the finite analytic erf profile, gives an
even shorter injectivity argument. The centered curve has speed
at most $c$, but its initial speed is strictly less than $c$ for
every nonzero $g$, since $\int F(1-F)|g|^2>0$.
Continuity makes its total length strictly less than $\pi$,
contradicting the antipodal endpoints. This observation supplies
no quantitative finite-device later-round yield.

\subsection{Complete thermal normalization after a fixed number of rounds}

Let $V_\tau=(3/2)h_\tau(B_\infty)Ff_1$ and
$w_m=K_\infty^{m-1}V_\tau$. The source is nonzero and injectivity
gives $w_m\ne0$ for every fixed $m$. Strong convergence of the
finite-box propagators, their uniform contraction bounds, and
the source convergence give
\begin{equation}
 K_{\tau,\ell}^{m-1}V_{0,\ell}(\tau)\longrightarrow w_m.
                                                        \label{eq:repeatlimit}
\end{equation}
The actual low columns after $m$ rounds are
\begin{equation}
 K_{\tau,\ell}^m u_E
 =b_E\ee^{-i\tau E}K_{\tau,\ell}^{m-1}V_{E,\ell}(\tau).
                                                        \label{eq:repeatcolumn}
\end{equation}
The scalar phase is $\ee^{-i\tau E}$, not a replacement of
$K^m$ by a commuting polynomial. Multiplication by a contraction
preserves the low Taylor approximation bounds and cannot increase
the singular values of the high-sector column operator. The
complete entropy bounds in Appendix~\ref{app:control} therefore
apply after normalization by the positive source coefficient,
giving
\begin{align}
 p_m&\sim\frac{\norm{w_m}^2}{m_1^2}\frac r s,
 \qquad\Tr\sigma_m^2\to1,\nonumber\\
 S(\sigma_m)&=O_{m,F}(s^{-2}\log s+s^{5/2}\ee^{-s\delta_F})\to0.
                                                        \label{eq:repeatentropy}
\end{align}
For the sharp wall one can take $\delta_F=1/2$ and $m_1=1$.
All fixed positive R\'enyi entropies also vanish. In particular
$p_{m+1}/p_m\to\norm{w_{m+1}}^2/\norm{w_m}^2>0$.
There is no lower bound uniform in $m$, no limit with $m$ growing,
and no assertion at every higher resonance.

\subsection{Two independent imperfect readouts}

With $\nu_j\in[0,1]$ and $\epsilon_j=1-\nu_j$, composition of the two maps in
Eq.~\eqref{eq:visibility} is exactly
\begin{align}
 \mathcal J_{\nu_2}(\mathcal J_{\nu_1}(\rho))
 ={}&\nu_1\nu_2K^2\rho(K^*)^2
 +\epsilon_1\nu_2K\mathcal D(\rho)K^*\nonumber\\
 &+\nu_1\epsilon_2\mathcal D(K\rho K^*)
 +\epsilon_1\epsilon_2\mathcal D^2(\rho).               \label{eq:repeatnoisemap}
\end{align}
All four terms are positive. If $p_1,p_2$ are the ideal probabilities,
the third branch before its scalar coefficient has trace $p_1/2$
and the fourth has trace $1/4$. For the second branch use
$U_B=U_A+2K$ and $U_Au_E=\ee^{-i\tau E}u_E$, giving
$KU_Bu_E=\ee^{-i\tau E}Ku_E+2K^2u_E$.
The $U_A$ term alone gives its lower trace $p_1/4$; the inequality
$\norm{a+2b}^2\le2\norm a^2+8\norm b^2$ gives the upper trace
$(3p_1+8p_2)/4$. Because both ideal probabilities are of order
$r/s$, Eq.~\eqref{eq:twonoise} makes the total false weight
negligible relative to $p_2$, proving source retention and ideal
leading yield. One effective visibility multiplying $K^2$ would
be a different channel.

The normalized singly false branches have uniformly bounded
entropy. For $K\mathcal D(\rho)K^*$ their low columns are
combinations of $Ku_E$ and $K^2u_E$, so two limiting source
vectors suffice, and their complete tails are bounded by sums
of the contracted analytic tails. Its trace is at least $p_1/4$,
so normalization is controlled. The other singly false branch
is a mixture of two unitary images of the one-round state.
Their weights relative to the useful branch are
$O(\epsilon_1+\epsilon_2)$, and their weighted entropies vanish.

The normalized all-false state $4\mathcal D^2(\rho)$ is a
mixture of four unitary images of $\rho$. Concavity and the
finite-label mixture bound give
\begin{equation}
 S(\rho)\le S(4\mathcal D^2(\rho))\le S(\rho)+\log4.
                                                        \label{eq:repeatfalseentropy}
\end{equation}
The partition and its first energy moment imply
$S(\rho)=\log N+O(1)$, $N=1/(r\sqrt s)$.
Consequently $\epsilon_1\epsilon_2\log N=o(r/s)$,
together with $\epsilon_1+\epsilon_2\to0$, pays the all-false
entropy and the finite mixture-label entropy, proving sufficiency.
Under Eq.~\eqref{eq:twonoise}, the all-false weight is
asymptotic to $\epsilon_1\epsilon_2/(4p_2)$; entropy concavity
and Eq.~\eqref{eq:repeatfalseentropy} prove necessity of this
strengthened product condition. These are results for the stated
independent visibility channels, not correlated dephasing or
unrecorded timing laws.

\section{A validated smooth-wall and noise calculation}
\label{app:certificate}

This calculation concerns the exact error-function wall in
Eq.~\eqref{eq:erf}. It does not replace that wall by a finite-dimensional
Hamiltonian. We describe the enclosures sufficiently to distinguish
analytic error control from numerical convergence.

\subsection{Resolving the transition with a controlled amplitude tail}

Put $w=1/40$, $e=\operatorname{erf}(5)$, $a_0=1-5w=7/8$. Define
a continuous plateau comparator by
\begin{equation}
 F_0(y)=\begin{cases}
 0,&y\le a_0,\\
 [\operatorname{erf}((y-1)/w)+e]/(2e),&a_0<y<2-a_0,\\
 1,&y\ge2-a_0.
 \end{cases}
\end{equation}
Monotonicity of the error function and direct subtraction inside
the middle interval give
\begin{equation}
 \norm{F_r-F_0}_\infty\le\varepsilon=(1-e)/2
                                      <7.688\times10^{-13}.
                                                               \label{eq:erfamp}
\end{equation}
This is an amplitude error. The mathematical validation collar
has length $1/4$, whereas the physical 10--90 width is
$2w\operatorname{erf}^{-1}(0.8)=0.9061938024\ldots r$.

For the spectral solve use $H_a=-\partial_y^2+aF_0$ with
$a=400,1600$, inverse temperature $8/400$, and time
$2\pi/1200$. Divide energies by 400 to return to
Eq.~\eqref{eq:erf}. The middle interval is divided into 50 cells
of length $h=1/200$. In a cell let $x\in[0,1]$ and
$z=-5+j/5$, so the error-function argument is $z+x/5$.
Taylor coefficients of $g(x)=\ee^{-(z+x/5)^2}$ follow from
$g'=-2(z+x/5)g/5$; integrating them supplies the coefficients
of $F_0$. The initial coefficient is enclosed with Arb's error
function. The solution coefficients obey the exact recurrence
\begin{equation}
 \begin{gathered}
 (n+1)(n+2)u_{n+2}
      \\=h^2\Bigl(a\sum_{j=0}^nF_j u_{n-j}
      \\-E u_n\Bigr).
                                                               \label{eq:erfrecurrence}
 \end{gathered}
\end{equation}
The value and derivative polynomials are retained through degree
48, with $u_{49}$ computed for the latter.

For an explicit common remainder, take the complex radius $R=4$
in the local variable. On that disk the error-function argument has
imaginary part at most $b=4/5$. Vertical integration from its real
part gives
\begin{equation}
 \begin{gathered}
 |\operatorname{erf}\zeta|\le1+2b\ee^{b^2}/\sqrt\pi,
 \qquad\\ |F_0|\le M_F^*=\frac{e+1+2b\ee^{b^2}/\sqrt\pi}{2e}.
 \end{gathered}
\end{equation}
For $Y=(u,u_x)$, the analytic first-order system has matrix
infinity norm at most $M=\max\{1,h^2(|E|+aM_F^*)\}$.
Gronwall along a radius and Cauchy's coefficient estimate imply
the uniform error bound for both retained components
\begin{equation}
 \eta\le\frac{\ee^{RM}\norm{Y(0)}_\infty}{(R-1)R^{48}}
                         \qquad(0\le x\le1).                 \label{eq:erfremainder}
\end{equation}
The physical derivative error is $\eta/h$.

To prevent interval wrapping, the propagated vector is recentered
at each cell endpoint, carrying a scalar Euclidean error in
$(u,u_y/k_*)$, with a fixed positive integer $k_*$ for that
eigenvalue bracket. The symmetric part of the real generator
has largest eigenvalue at most
\begin{equation}
             \frac{|k_*^2-E|+a}{2k_*}.
\end{equation}
An incoming error is therefore multiplied by at most the
exponential of $h$ times this bound. Local endpoint ball radii
and Eq.~\eqref{eq:erfremainder} are added outwardly. The same
growth bound controls every interior point used in integrals.
Energy uncertainty is retained in each recurrence and exponent;
recentering discards no error term.

\subsection{Root labels, overlaps, and the full accepted Gram matrix}

The left plateau solution is $\sin(\sqrt E\,y)$ and the right
solution is boundary-adapted sine or hyperbolic sine. At the right
collar endpoint the matching determinant is
\begin{equation}
 \begin{cases}
 u_b\cos(qa_0)+v_b\sin(qa_0)/q,&E>a,\quad q=\sqrt{E-a},\\
 u_b+v_b\tanh(\kappa a_0)/\kappa,&E<a,\quad\kappa=\sqrt{a-E}.
 \end{cases}                                                  \label{eq:erfroot}
\end{equation}
The threshold is its continuous limit. Floating integration only
proposes rational root brackets. The analytic ODE recomputes their
endpoint signs. The free plateau nodes are counted exactly; on
each collar cell, $u''+(E-aF_0)u=0$ and $E-aF_0\le E$ imply
that two zeros are separated by at least $\pi/\sqrt E$.
The checker verifies that the cell is shorter than this distance
and that its endpoint values are nonzero. Opposite signs then
mean exactly one zero, equal signs none. Consecutive Sturm indices,
ordered disjoint brackets, and the first omitted neighbor certify
the spectral prefix without missing an intermediate level.

Cell overlap integrals multiply the enclosed polynomials. If their
supremum bounds are $U,V$ and errors $\eta,\zeta$, the missing
integral is at most $h(\eta V+\zeta U+\eta\zeta)$. Exact signed
sine and hyperbolic-sine integrals supply both plateaus. This
constructs normalization and signed overlaps for 30 initial and
126 final retained modes. The first omitted final eigenvalue
and the overlap identity bound every omitted final overlap mass
$R_n$ by Eq.~\eqref{eq:finalcounttail}, replacing $Q$ by $F_0$.
In the $H_a$ energy units the numerator is
$1200^2\norm{F_0u_n}^2$; in dimensionless energies, after division
by 400, it is $9\norm{F_0u_n}^2$. The numerator norm is bounded
by the full collar and right-plateau masses.
Equation~\eqref{eq:initialcounttail} bounds
every omitted initial Gibbs weight.

There is a useful factor of two in the Gram tail. For real
eigenfunctions let
$V_{nm}(t)=\sum_j O_{jn}O_{jm}\ee^{-it\lambda_j}$ in dimensionless
energies. The exact initial-basis Gram matrix is
\begin{equation}
 \begin{gathered}
 G_{nm}=\ip{K_tu_n}{K_tu_m}
 \\=\frac{\delta_{nm}}2-\frac14\left[
       \ee^{itE_n}V_{nm}+\overline{\ee^{itE_m}V_{nm}}\right].
                                                               \label{eq:erfgram}
 \end{gathered}
\end{equation}
The missing part of $V_{nm}$ is at most $\sqrt{R_nR_m}$,
so the Gram error is at most $\sqrt{R_nR_m}/2$.
Let $a_n=\ee^{-sE_n}$ and $T=\sum_na_nG_{nn}$. Positivity gives
\begin{equation}
 p=T/Z_1,\qquad
 \Tr\sigma^2=\frac{\sum_{n,m}a_na_m|G_{nm}|^2}{T^2}.          \label{eq:erfpurity}
\end{equation}
For a lower purity bound the retained $(n,m)$ entry contributes
$[\max\{0,\underline{|G_{nm}^{\rm prefix}|}
-\sqrt{R_nR_m}/2\}]^2$, where the underline denotes its interval
lower bound. All missing squared entries are nonnegative.
Upper-bound $T$ by the retained diagonal
upper sum plus the full omitted initial trace, since $G_{nn}\le1$.
Lower and upper partition bounds give the acceptance interval.
Input purity is $Z_1(2s)/Z_1(s)^2$, with both heat tails paid.

\subsection{Passing to the actual analytic wall}

Let $S=\ee^{-sA/2}/\sqrt{Z_1}$ and
$S_0=\ee^{-sA_0/2}/\sqrt{Z_{1,0}}$ be Hilbert--Schmidt
purifications. Min--max and Eq.~\eqref{eq:erfamp} give
$\ee^{-s\varepsilon}Z_{1,0}\le Z_1\le
\ee^{s\varepsilon}Z_{1,0}$. Duhamel at heat time $t=s/2$,
with Schatten H\"older, bounds the integrand by
\begin{equation}
 \begin{gathered}
 \norm{\ee^{-(t-u)A}(F_r-F_0)\ee^{-uA_0}}_2
 \\
 \le\varepsilon Z_1^{(t-u)/(2t)}Z_{1,0}^{u/(2t)}
 \\
 \le\varepsilon\ee^{(t-u)\varepsilon}\sqrt{Z_{1,0}}.
 \end{gathered}
\end{equation}
After integration and normalization of the two vectors this gives
\begin{equation}
       \norm{S-S_0}_2\le d_{\rm prep}:=s\varepsilon\ee^{s\varepsilon/2}.
                                                               \label{eq:erfprep}
\end{equation}
Unitary Duhamel also gives
$\norm{K_\tau(F_r)-K_\tau(F_0)}\le(5/2)|\tau|\varepsilon$.
Thus the unnormalized output purifications differ by at most
$h=d_{\rm prep}+(5/2)|\tau|\varepsilon$.

More generally, if a reference output purification has squared norm
in $[p_-,p_+]$ and its state has purity at least $P_-$, a vector
error at most $h<\sqrt{p_-}$ yields
\begin{equation}
 (\sqrt{p_-}-h)^2\le p\le(\sqrt{p_+}+h)^2,
 \qquad \Tr\sigma^2\ge P_- -\frac{2h}{\sqrt{p_-}}.            \label{eq:purificationbound}
\end{equation}
The probability bounds are the norm triangle inequalities. For the
state, project the reference vector onto the line through the new
one; the sine of their angle is at most $h/\sqrt{p_-}$.
The pure-state trace-distance identity and partial-trace
contractivity give the displayed purity loss, because purity is
1-Lipschitz in trace norm on density operators:
\begin{equation}
 |\Tr(\sigma^2-\rho^2)|
 =|\Tr[(\sigma-\rho)(\sigma+\rho-\id)]|
 \le\norm{\sigma-\rho}_1.
\end{equation}
The last inequality uses $-\id\le\sigma+\rho-\id\le\id$.
Input purity changes by at most $2d_{\rm prep}$.
These bounds give the actual analytic-wall nominal acceptance
$0.0025635744408<p_0<0.0025635744431$, nominal purity
$P_0>0.9664657000$, and input purity below $0.386060$.

\subsection{Unrecorded time distributions and recorded outcomes}

Since $A$ commutes with its Gibbs square root,
$B=A+3F_r$ and unitarity give Eq.~\eqref{eq:jitter} by
differentiating in Hilbert--Schmidt norm. In particular
$\norm{(K_t-K_\tau)S}_2\le L|t-\tau|$.
For any independent duration law $\mu$ supported on
$[\tau-\Delta,\tau+\Delta]$, regard $K_tS$ as a vector in
$L^2(\mu;\mathfrak S_2)$. Its distance from the constant
reference $K_\tau S$ is at most $L\Delta$. Partial trace over
this duration register gives exactly the unrecorded mixture
$\overline R=\int K_t\rho K_t^*\mu(\dd t)$.
Equation~\eqref{eq:purificationbound} therefore applies to this
mixture, not just separately to states in duration-dependent frames.

Both moments in $L$ are bounded for the actual $F_r$. For retained
comparator energies $E_n$ in dimensionless units, the actual
unnormalized second-energy contribution is bounded by
$(E_n+\varepsilon)^2\ee^{-s(E_n-\varepsilon)}$.
For omitted modes, put
$c=sr^2\pi^2/4$, $x_0=r^2\pi^2(n_0+1)^2/4$ and
$q_0=\ee^{-c(2n_0+3)}$. The checker verifies $x_0>2/s$.
Monotonicity of $x^2\ee^{-sx}$ above $x_0$ and
$(n_0+1+j)^4\le(n_0+1)^4(j+1)^4$ give the full tail
\begin{equation}
 x_0^2\ee^{-c(n_0+1)^2}
              \frac{1+11q_0+11q_0^2+q_0^3}{(1-q_0)^5}.
                                                               \label{eq:momenttail}
\end{equation}
Divide by the actual partition lower bound
$\ee^{-s\varepsilon}Z_{1,0,-}$. For $\Tr\rho_0F_0^2$, bound
$F_0^2$ by its right endpoint in each validation cell and multiply
by the certified cell norm. Add the entire right-plateau norm,
where $F_0=1$, and the complete omitted Gibbs tail.
Then
\begin{equation}
 \Tr\rho F_r^2\le\Tr\rho_0F_0^2+2d_{\rm prep}+2\varepsilon.
\end{equation}
The resulting bounds are
$\Tr\rho A^2<0.015989357$, $\Tr\rho F_r^2<0.003950444$,
and $L<0.220727856$.

To include visibility and detector assignments, set
\begin{equation}
 \overline D=\int\frac{U_A\rho U_A^*+U_B\rho U_B^*}{2}\mu(\dd t),
 \qquad\Tr\overline D=1.
\end{equation}
Both branches use the same duration law. With false-positive
probability $e$ and false-negative probability $d$, the recorded
dark operator is exactly
\begin{equation}
 \begin{gathered}
 R_{\rm obs}=a\overline R+b\overline D,\qquad\\
 a=(1-d-e)\nu,\quad b=(1-d-e)(1-\nu)/2+e.                     \label{eq:recordedmap}
 \end{gathered}
\end{equation}
This follows from $(1-d)\mathcal J_\nu+e\mathcal J_{\rm bright}$
and $\mathcal J_\nu+\mathcal J_{\rm bright}=\overline D$.
For $d+e\le1$ both coefficients are nonnegative. Positivity
of the Hilbert--Schmidt cross term gives
\begin{equation}
 \Tr\sigma_{\rm obs}^2\ge
 \left(\frac{a\overline p}{a\overline p+b}\right)^2
                  \Tr\overline\sigma^2,
 \qquad p_{\rm obs}=a\overline p+b.                          \label{eq:recordedpurity}
\end{equation}
Using $\Delta=0.001$, $\nu=0.9999$, $e=10^{-5}$, $d=0.01$
and outward endpoints in these monotone formulas gives
Eq.~\eqref{eq:certificate}. All probabilities are conditional
on the specified physical channels; generic correlated preparation
or detection errors are different models.

The accompanying scripts enclose every root, ODE remainder, overlap,
Gram entry, Gibbs tail, and scalar propagation at 256-bit ball
precision. Saved source hashes identify the scripts actually
replayed. Their finite prefixes are computational devices inside
the inequalities above and never project the physical output.

\section{Finite spectral enclosures and statistical reconstruction}
\label{app:counts}

This appendix distinguishes deterministic continuum error from
statistical counting error. Both are needed to use the source
reconstruction in an actual finite experiment. Throughout this
appendix the device is the sharp-wall Hamiltonian in
Eq.~\eqref{eq:ham}; the noisy smooth example is separate.

\subsection{Complete finite-device probabilities}

For each initial eigenstate, completeness and $B-A=3Q$ give
\begin{align}
 \sum_j|O_{jn}|^2&=1,\nonumber\\
 (\lambda_j-E_n)O_{jn}&=3\ip{\phi_j}{Qu_n},\label{eq:countoverlap}\\
 \sum_j(\lambda_j-E_n)^2|O_{jn}|^2&=9\norm{Qu_n}^2,
 \qquad\nonumber\\ O_{jn}&=\ip{\phi_j}{u_n}.\nonumber
\end{align}
Take a complete initial prefix of length $n_0$ and a complete final
prefix of length $j_0$, including certified lower endpoints for the
first omitted levels. If $\lambda_{j_0+1}\ge L_*>E_n$, the missing
final overlap mass obeys
\begin{equation}
 \begin{gathered}
 R_n:=\sum_{j>j_0}|O_{jn}|^2
 \\
 \le\min\left\{1-\sum_{j\le j_0}|O_{jn}|^2,
                \frac{9\norm{Qu_n}^2}{(L_*-E_n)^2}\right\}.
                                                               \label{eq:finalcounttail}
 \end{gathered}
\end{equation}
The upper interval endpoint of each expression is used. For the full
omitted initial sector, min--max comparison with the free Dirichlet
interval gives
\begin{equation}
 T_{n_0}:=\sum_{n>n_0}\ee^{-sE_n}
 \le\frac{\ee^{-b(n_0+1)^2}}
           {1-\ee^{-b(2n_0+3)}},\qquad b=s r^2\pi^2/4.
                                                               \label{eq:initialcounttail}
\end{equation}
Indeed the successive exponents beyond $n_0+1$ increase by at least
$b(2n_0+3)$. This tail includes every omitted state above the barrier;
retained states in that sector remain in the finite sum.

Let $Z_-$ and $Z_+$ enclose
$\sum_{n\le n_0}\ee^{-sE_n}+[0,T_{n_0}]$, and define the finite
positive numerator
\begin{equation}
 P_{n_0j_0}(t)=\sum_{n\le n_0,j\le j_0}\ee^{-sE_n}|O_{jn}|^2
                         \sin^2[t(\lambda_j-E_n)/2].
\end{equation}
Equations~\eqref{eq:ramseytpm} and \eqref{eq:finalcounttail} then give
\begin{equation}
 \frac{P_{n_0j_0}(t)}{Z_+}\le p_{r,s}(t)
 \le\frac{P_{n_0j_0}(t)+\sum_{n\le n_0}\ee^{-sE_n}R_n+T_{n_0}}
              {Z_-}.                                         \label{eq:countenclosure}
\end{equation}
The finite sums themselves are evaluated with outward interval
arithmetic, so their lower and upper endpoints are used on the
corresponding sides. Neither tail uses oscillatory cancellation.

For clarity, all eigenvalues can be enclosed without a spatial
discretization. For a wall of height $h$ on the right half of $(0,2)$,
write $k=\sqrt E/r$, $q=\sqrt{E-h}/r$. Matching $\sin(ky)/k$ with
the right boundary-adapted function $\sin[q(2-y)]/q$ gives the
entirely continued determinant
\begin{equation}
       \cos k\,\frac{\sin q}{q}+\frac{\sin k}{k}\cos q=0.      \label{eq:countdet}
\end{equation}
Imaginary $q$ gives the hyperbolic formula; both sine quotients are
continued with value one at zero. Rational sign brackets, disjoint in increasing order, and
certified Sturm node counts prove that the enclosed roots have
consecutive indices starting at one. Analytic sine and hyperbolic-sine
integrals give their normalization and signed cross overlaps.
Floating-point root searches only propose brackets; signs and
indices are rechecked with ball arithmetic \cite{JohanssonArb}.

The data in Sec.~\ref{sec:counts} use $n_0=64$ and $j_0=1260$ at
$(r,s)=(1/200,128)$. The weighted, scaled Gibbs-tail allowance is
below $5.612\times10^{-11}$; the final spectral-tail allowance is
below $1.297776\times10^{-5}$. These errors are already included
in Eq.~\eqref{eq:countenclosure}, not appended a second time.

\subsection{The exact source quadrature and its tails}

For $c=1$, use $T=32$, the 256 nodes $t_i=(2i+1)/16$,
$i=0,\ldots,255$, and the positive weights in Eq.~\eqref{eq:weights}.
Their exact sum is
\begin{equation}
 \begin{gathered}
 W=\frac4\pi\left[\operatorname{Si}(32)
                       -\frac{1-\cos32}{32}\right],
 \qquad\\ W=1.95959365427766\ldots .                            \label{eq:weighttotal}
 \end{gathered}
\end{equation}
The weight is fixed by the reconstruction identity, not fitted to
the value of $\Phi(1)$.

For the low-energy part of $C(t)$, substitute $E=4\sin^2\theta$ in
Eq.~\eqref{eq:sharpdensity}. Its integrand on $[0,\pi/2]$ becomes
\begin{equation}
 \frac{9t^2}{\pi}
 \frac{\sin^2\theta\cos\theta}{(1+2\cos\theta)^2}
                 \sinc^2(2t\sin^2\theta),                   \label{eq:sourcequadlow}
\end{equation}
where $\sinc z=\sin z/z$ is continued at zero. Above four, set
$E=4+q^2$, $k=\sqrt{4+q^2}$. The measure density in $q$ is
\begin{equation}
 h(q)=\frac{18q(1+kq)}
          {\pi(q^2+4)^2(k+q)(q^2+1)^2}.                       \label{eq:sourcequadhigh}
\end{equation}
Integrate the actual sine-square response over $0\le q\le4$.
The complete omitted mass is enclosed by the nonsingular finite
integral, obtained from $q=1/u$,
\begin{equation}
 \begin{gathered}
 \nu((20,\infty))=
 \\
 \int_0^{1/4}\frac{18u^4(u^2+\sqrt{1+4u^2})}
 {\pi(1+4u^2)^2(1+\sqrt{1+4u^2})(1+u^2)^2}\,\dd u.
                                                               \label{eq:sourcetail}
 \end{gathered}
\end{equation}
It lies within $0.0004026117\pm4.36\times10^{-11}$.
For every time node its missing sine-square contribution belongs to
$[0,\nu((20,\infty))]$.

For $\Phi(1)$ the low-angle integral is split at $\theta=\pi/6$.
Before and after this point the integrands are respectively
\begin{equation}
 \frac{9\cos\theta}{\pi(1+2\cos\theta)^2},\qquad
 \frac{9\cos\theta}{4\pi\sin^2\theta(1+2\cos\theta)^2}.
\end{equation}
The complete mass above four from Eqs.~\eqref{eq:sourcequadhigh}
and \eqref{eq:sourcetail} is then added. Analytic ball quadrature
\cite{JohanssonIntegration} encloses all finite integrals, producing
Eq.~\eqref{eq:phinumeric} without a missing high-energy tail.

Writing $S_C=\sum_iw_iC(t_i)$, the certificate keeps separate
enclosures for $S_C-\Phi(1)$, $S-S_C$, and
$\sum_iw_i|\lambda p(t_i)-C(t_i)|$. At the stated device,
\begin{align}
 |S-\Phi(1)|&<0.000479010749,\nonumber\\
 |S_C-\Phi(1)|&<0.047699256289,\label{eq:systematiccounts}\\
 \sum_iw_i|\lambda p(t_i)-C(t_i)|&<0.047233223348.\nonumber
\end{align}
The small direct error includes cancellation between two biases of
opposite sign. The separate absolute budget is below
$0.094932479636$. Because both the full source target and the finite
source grid are enclosed, this already pays the source time-window
and mesh errors; a further time-tail allowance would double-count
them. These numerical statements concern this known sharp source,
not a uniformly stable inverse of an unspecified potential.

\subsection{Exact count intervals and the preparation budget}

Assume independent trials with the same node probabilities $p(t_i)$.
Let $X$ be the dark-event indicator after randomly choosing a node
with probability $w_i/W$. Since this choice is made independently
on every trial, $X$ is Bernoulli with
$\bar p=\sum_i(w_i/W)p(t_i)$; a mixture of Bernoulli indicators is
still Bernoulli. Therefore $k=\sum_{j=1}^nX_j$ is binomial.
One exact equal-tailed interval \cite{ClopperPearson} is
\begin{align}
 L_n(k)&=\begin{cases}
 0,&k=0,\\
 \operatorname{BetaQuantile}(\alpha/2;k,n-k+1),&k>0,
 \end{cases}\nonumber\\
 U_n(k)&=\begin{cases}
 1,&k=n,\\
 \operatorname{BetaQuantile}(1-\alpha/2;k+1,n-k),&k<n.
 \end{cases}                                                  \label{eq:cp}
\end{align}
Each endpoint inverts one binomial tail of probability $\alpha/2$;
the union bound gives coverage at least $1-\alpha$. In particular
$U_n(0)=1-(\alpha/2)^{1/n}$, so zero successes do not require a
Gaussian approximation. Combining the interval with deterministic
$d_\pm$ proves Eq.~\eqref{eq:confidence}.

For the predeclared budget, a centered Bernoulli variable $Z$ has
$\mathbb E|Z|^j\le\Var Z$ for every integer $j\ge2$. Expanding
the exponential and using $j!\ge2\,3^{j-2}$ gives
\begin{equation}
 \log\mathbb E\ee^{uZ}\le
       \frac{u^2\Var Z}{2(1-|u|/3)},\qquad |u|<3.
\end{equation}
Exponential Markov for the two signs and independence then yield
\begin{equation}
 \mathbb P(|\widehat S-S|\ge\epsilon)
 \le2\exp\left[-\frac{n\epsilon^2}
                   {2\lambda W(D+\epsilon/3)}\right],
                                                               \label{eq:bernstein}
\end{equation}
because $\lambda W\bar p=S\le D$.
This proves Eq.~\eqref{eq:shots}. A positive upper bound on the
\emph{actual} scaled response is required: subtracting a known
background mean would not remove its Bernoulli variance.

The implementation uses exact dyadic weights
$\widetilde w_i=2^{-80}\lfloor2^{80}w_i\rfloor$, with each integer
numerator certified uniquely. If $p_i\le U_i/\lambda$ and
$|w_i-\widetilde w_i|\le e_i$, the added bias is at most
\begin{equation}
             B_w=\sum_i e_iU_i.                              \label{eq:roundweight}
\end{equation}
Run the interval and budget formulas with $\widetilde W$ and the
exact dyadic sampling law, widening the bias endpoints by $B_w$.
Here $B_w<5.35\times10^{-23}$. A sampling-law implementation with
total-variation error $\zeta$ would require the further allowance
$\lambda\widetilde W\zeta$; the stated budgets specify the exact
discrete law.

For $\alpha=0.05$, use $\epsilon=0.05$ with the first line of
Eq.~\eqref{eq:systematiccounts}, or $\epsilon=0.005$ with its two
separate absolute errors. Outward rounding of Eq.~\eqref{eq:shots}
gives the counts displayed in Sec.~\ref{sec:counts}, and respective
total error bounds $0.050479011$ and $0.099932480$.
No measured count is fabricated by this design calculation.

\subsection{Allocation and duration are part of the resource model}

The pooled protocol is simple but not universally variance optimal.
If $v_i=p_i(1-p_i)$ and fixed counts $n_i=nq_i$ are used at the
nodes, the idealized unrounded variance is
\begin{equation}
 \frac{\lambda^2}{n}\sum_i\frac{w_i^2v_i}{q_i}.
\end{equation}
Cauchy--Schwarz minimizes it at
$q_i\propto w_i\sqrt{v_i}$ on nonzero-variance nodes. This is an
oracle allocation; learning the $p_i$ requires its own data. Equal
shots have between one and $m$ times that minimal variance for
$m$ nodes, provided at least one $v_i>0$. Randomizing with
$q_i=w_i/W$ adds, relative to the same
fixed allocation, the nonnegative variance
$\lambda^2W^2\Var_q(p_i)/n$.

A usable confidence-bound design uses certified $U_i$ instead.
For any positive distribution $q$, one trial has value
$Y=\lambda w_IX/q_I$, expectation $S$, variance at most
$\lambda A(q)$, and range at most $\lambda B(q)$, where
\begin{equation}
 A(q)=\sum_iw_i^2U_i/q_i,\qquad B(q)=\max_iw_i/q_i.
\end{equation}
The same argument as Eq.~\eqref{eq:bernstein} replaces $W(D+\epsilon/3)$
by $A(q)+\epsilon B(q)/3$. Its minimization is a finite convex
problem. If all $U_i=0$, then $S=0$ and no sampling is required.
Otherwise, at fixed $b\ge W$ and constraints $q_i\ge w_i/b$, the
minimizer has
\begin{equation}
 q_i(b)=\max\{w_i/b,\ w_i\sqrt{U_i}/t_b\},\qquad
                \sum_iq_i(b)=1,                              \label{eq:waterfill}
\end{equation}
with the limiting convention when a set of $U_i$ vanishes.
Minimizing $A(q(b))+\epsilon b/3$ over $b\ge W$ gives the best
bound of this form, not an optimal quantum sensing theorem.
The displayed numerical budgets use the simpler pooled design.

Each trial consumes one complete Gibbs preparation, a fresh pure
ancilla, a controlled pair of branches, and a readout. At the
256-node design the mean dimensionless duration is
$2.616668820164279\ldots$, and its maximum is $31.9375$.
For any fixed $n$, the mean total controlled duration is
\begin{equation}
 \mathbb E\mathcal T=n\frac{\hbar r^2}{E_0}\sum_iq_it_i,
 \qquad\mathcal T\le n\frac{\hbar r^2T}{E_0}.
\end{equation}
Preparation, reset, switching, and readout overheads are additional.
The probability of retaining a state is not a license to omit
failed trials from either the confidence interval or this cost.

\begin{thebibliography}{99}

\bibitem{Dorner2013}
R.~Dorner, S.~R.~Clark, L.~Heaney, R.~Fazio, J.~Goold, and V.~Vedral,
Extracting quantum work statistics and fluctuation theorems by
single-qubit interferometry,
\href{https://doi.org/10.1103/PhysRevLett.110.230601}
{Phys. Rev. Lett. \textbf{110}, 230601 (2013)}.

\bibitem{Mazzola2013}
L.~Mazzola, G.~De Chiara, and M.~Paternostro,
Measuring the characteristic function of the work distribution,
\href{https://doi.org/10.1103/PhysRevLett.110.230602}
{Phys. Rev. Lett. \textbf{110}, 230602 (2013)}.

\bibitem{Batalhao2014}
T.~B.~Batalh\~ao, A.~M.~Souza, L.~Mazzola, R.~Auccaise,
R.~S.~Sarthour, I.~S.~Oliveira, J.~Goold, G.~De Chiara,
M.~Paternostro, and R.~M.~Serra,
Experimental reconstruction of work distribution and study of
fluctuation relations in a closed quantum system,
\href{https://doi.org/10.1103/PhysRevLett.113.140601}
{Phys. Rev. Lett. \textbf{113}, 140601 (2014)}.

\bibitem{Nakazato2003}
H.~Nakazato, T.~Takazawa, and K.~Yuasa,
Purification through Zeno-like measurements,
\href{https://doi.org/10.1103/PhysRevLett.90.060401}
{Phys. Rev. Lett. \textbf{90}, 060401 (2003)}.

\bibitem{Oyarce2019}
F.~Oyarce and M.~Orszag,
Photon cooling by dispersive atom-field coupling with atomic postselection,
\href{https://arxiv.org/abs/1904.03839}{arXiv:1904.03839 (2019)}.

\bibitem{Pushpan2025}
C.~B.~Pushpan, T.~K.~Konar, A.~Sen(De), and A.~K.~Pal,
Developing universal logical state-purification strategy for quantum
error correcting codes,
\href{https://arxiv.org/abs/2502.01393}{arXiv:2502.01393 (2025)}.

\bibitem{Francica2019}
G.~Francica, J.~Goold, and F.~Plastina,
Role of coherence in the nonequilibrium thermodynamics of quantum systems,
\href{https://doi.org/10.1103/PhysRevE.99.042105}
{Phys. Rev. E \textbf{99}, 042105 (2019)}.

\bibitem{Varizi2022}
A.~D.~Varizi, R.~C.~Drumond, and G.~T.~Landi,
Quantum quench thermodynamics at high temperatures,
\href{https://doi.org/10.1103/PhysRevA.105.062218}
{Phys. Rev. A \textbf{105}, 062218 (2022)}.

\bibitem{Smacchia2012}
P.~Smacchia and A.~Silva,
Universal energy distribution of quasiparticles emitted in a local
time-dependent quench,
\href{https://doi.org/10.1103/PhysRevLett.109.037202}
{Phys. Rev. Lett. \textbf{109}, 037202 (2012)}.

\bibitem{Knap2012}
M.~Knap, A.~Shashi, Y.~Nishida, A.~Imambekov, D.~A.~Abanin,
and E.~Demler,
Time-dependent impurity in ultracold fermions: Orthogonality catastrophe
and beyond,
\href{https://doi.org/10.1103/PhysRevX.2.041020}
{Phys. Rev. X \textbf{2}, 041020 (2012)}.

\bibitem{Kuttler2014}
H.~K\"uttler, P.~Otte, and W.~Spitzer,
Anderson's orthogonality catastrophe for one-dimensional systems,
\href{https://doi.org/10.1007/s00023-013-0287-z}
{Ann. Henri Poincar\'e \textbf{15}, 1655--1696 (2014)}.
\bibitem{Sindona2013}
A.~Sindona, N.~Lo Gullo, J.~Goold, and F.~Plastina,
Statistics of the work distribution for a quenched Fermi gas,
\href{https://arxiv.org/abs/1309.2669}{arXiv:1309.2669 (2013)}.

\bibitem{Guarnieri2024}
G.~Guarnieri, J.~Eisert, and H.~J.~D.~Miller,
Generalized linear response theory for the full quantum work statistics,
\href{https://doi.org/10.1103/PhysRevLett.133.070405}
{Phys. Rev. Lett. \textbf{133}, 070405 (2024)}.

\bibitem{Ahmad2022}
S.~Ali Ahmad and A.~R.~H.~Smith,
Finite-resolution ancilla-assisted measurements of quantum work distributions,
\href{https://doi.org/10.1103/PhysRevA.105.052411}
{Phys. Rev. A \textbf{105}, 052411 (2022)}.

\bibitem{Alonso2026}
D.~Alonso and A.~Ruiz-Garc\'ia,
Operational relation between one-time and two-time work protocols and
measurement-resolution-induced transitions in quantum work statistics,
\href{https://arxiv.org/abs/2607.16951}{arXiv:2607.16951 (2026)}.

\bibitem{Hegerfeldt1998}
G.~C.~Hegerfeldt,
Instantaneous spreading and Einstein causality in quantum theory,
\href{https://doi.org/10.1002/andp.199851007-817}
{Ann. Phys. (Leipzig) \textbf{7}, 716--725 (1998)};
\href{https://arxiv.org/abs/quant-ph/9809030}{arXiv:quant-ph/9809030}.

\bibitem{ChandruRareWork}
R.~Arun Chandru,
Rare work costs under bounded-action quantum confinement control,
\href{https://doi.org/10.5281/zenodo.23180314}
{Zenodo preprint, doi:10.5281/zenodo.23180314 (2026)}.

\bibitem{JohanssonArb}
F.~Johansson,
Arb: Efficient arbitrary-precision midpoint-radius interval arithmetic,
\href{https://arxiv.org/abs/1611.02831}{arXiv:1611.02831 (2016)}.

\bibitem{JohanssonIntegration}
F.~Johansson,
Numerical integration in arbitrary-precision ball arithmetic,
\href{https://arxiv.org/abs/1802.07942}{arXiv:1802.07942 (2018)}.

\bibitem{ClopperPearson}
C.~J.~Clopper and E.~S.~Pearson,
The use of confidence or fiducial limits illustrated in the case of the binomial,
\href{https://doi.org/10.1093/biomet/26.4.404}
{Biometrika \textbf{26}, 404--413 (1934)}.

\end{thebibliography}

\end{document}
